\documentclass[12pt, a4paper]{article}
\usepackage[T2A]{fontenc}
\usepackage[utf8]{inputenc}
\usepackage[russian]{babel}
\usepackage{amsmath, amssymb, amsthm}
\usepackage{geometry}
\usepackage{hyperref}
\usepackage{enumitem}

\geometry{left=2.5cm, right=2.5cm, top=2.5cm, bottom=2.5cm}

\newtheorem{theorem}{Теорема}
\newtheorem{lemma}[theorem]{Лемма}
\newtheorem{corollary}[theorem]{Следствие}
\newtheorem{proposition}{Предложение}
\theoremstyle{definition}
\newtheorem{definition}{Определение}
\newtheorem{example}{Пример}
\theoremstyle{remark}
\newtheorem*{remark}{Замечание}

\newcommand{\Z}{\mathbb{Z}}
\newcommand{\Q}{\mathbb{Q}}
\newcommand{\R}{\mathbb{R}}
\renewcommand{\C}{\mathbb{C}}
\newcommand{\F}{\mathbb{F}}
\newcommand{\Ker}{\mathrm{Ker}\,}
\renewcommand{\Im}{\mathrm{Im}\\,}

\title{\textbf{Ответы на вопросы к экзамену по алгебре}\\\large 1-64, Весна 2026}
\author{}
\date{}

\begin{document}
	\maketitle
	\tableofcontents
	\newpage

%──────────────────────────────────────────────────────────────────────────────
\section{Кольцо. Делители нуля, регулярные и обратимые элементы. Область целостности}
\addcontentsline{toc}{section}{Вопрос 1. Кольцо. Делители нуля, регулярные и обратимые элементы. Область целостности}

% Вопрос 1. Определение кольца. Делители нуля и регулярные элементы. Доказать, обратимый элемент является регулярным...

\textbf{Определение.} Кольцо $\langle R, +, \cdot \rangle$ — это множество с двумя бинарными операциями, являющееся абелевой группой по сложению, полугруппой по умножению, в котором выполняется дистрибутивность умножения относительно сложения.

\textbf{Делители нуля и регулярные элементы.} Ненулевой элемент $a \in R$ называется левым (правым) делителем нуля, если существует ненулевой $b \in R$, такой что $ab = 0$ ($ba = 0$). Элемент, не являющийся делителем нуля (ни левым, ни правым) и отличный от нуля, называется регулярным.

\textbf{Теорема.} В кольце с единицей обратимый элемент является регулярным.
\textbf{Доказательство.} Пусть $a \in R$ — обратимый элемент (т.е. $\exists a^{-1}: a^{-1}a = aa^{-1} = 1$). Предположим от противного, что $a$ — делитель нуля. Тогда существует $b \neq 0$ такое, что $ab = 0$. Умножим обе части слева на $a^{-1}$:
$$ a^{-1}(ab) = a^{-1} \cdot 0 \implies (a^{-1}a)b = 0 \implies 1 \cdot b = 0 \implies b = 0. $$
Получили противоречие с тем, что $b \neq 0$. Следовательно, $a$ не может быть делителем нуля.

\textbf{Группа обратимых элементов.} 
\textbf{Теорема:} В кольце $R$ с единицей множество всех обратимых элементов $R^*$ образует группу по умножению.
\textbf{Доказательство:}
1. Замкнутость: Если $a, b \in R^*$, то $(ab)^{-1} = b^{-1}a^{-1} \in R$, значит $ab \in R^*$.
2. Ассоциативность: Наследуется из кольца.
3. Нейтральный элемент: Единица кольца $1 \in R^*$, так как $1 \cdot 1 = 1$.
4. Обратный элемент: Если $a \in R^*$, то $a^{-1}$ тоже обратим (его обратный — это $a$), значит $a^{-1} \in R^*$. Ч.т.д.

\textbf{Область целостности} — это коммутативное кольцо с единицей без нетривиальных делителей нуля ($ab=0 \implies a=0$ или $b=0$).

%──────────────────────────────────────────────────────────────────────────────
\section{Поле. Подкольца, унитарные подкольца и подполя}
\addcontentsline{toc}{section}{Вопрос 2. Поле. Подкольца, унитарные подкольца и подполя}

% Вопрос 2. Определение поля... Докажите характеристическое свойство подполя. Примеры подколец...

\textbf{Определение поля.} Поле $F$ — это коммутативное кольцо с единицей ($1 \neq 0$), в котором каждый ненулевой элемент обратим.
\textbf{Мультипликативная группа поля} $F^*$ — это множество всех ненулевых элементов поля (образует абелеву группу по умножению). Регулярные элементы поля — все элементы, кроме нуля.
\textbf{Свойство:} Поле является областью целостности, так как любой ненулевой (а значит, обратимый) элемент регулярный (доказано в Вопросе 1).

\textbf{Определения:} Подкольцо $S \subset R$ — подмножество, само являющееся кольцом относительно тех же операций. Унитарное подкольцо — подкольцо, содержащее единицу надкольца. Подполе — подмножество поля, само являющееся полем.

\textbf{Характеристическое свойство подполя.}
\textbf{Критерий:} Подмножество $L$ поля $F$ является подполем тогда и только тогда, когда:
1) $L$ содержит не менее двух элементов;
2) $\forall a, b \in L \implies a - b \in L$;
3) $\forall a \in L, \forall b \in L \setminus \{0\} \implies ab^{-1} \in L$.

\textbf{Доказательство:} 
Необходимость очевидна из определения поля.
Достаточность: 
Из (2) при $a=b$ следует $0 = a-a \in L$. 
При $a=0$ имеем $0-b = -b \in L$, то есть есть противоположные. Вкупе с замкнутостью по вычитанию это дает подгруппу по сложению.
Так как $|L| \geq 2$, существует $b \in L \setminus \{0\}$. Из (3) при $a=b$ имеем $b b^{-1} = 1 \in L$.
Из (3) при $a=1$ получаем $1 b^{-1} = b^{-1} \in L$ (существование обратного).
Наконец, для умножения: так как $b \in L \implies b^{-1} \in L$, то подставив в (3) $b^{-1}$ вместо $b$, получим $a (b^{-1})^{-1} = ab \in L$ (замкнутость по умножению). Ч.т.д.

\textbf{Примеры подколец с единицами:}
1. \textit{Совпадают:} Подкольцо целых чисел $\mathbb{Z}$ в поле рациональных $\mathbb{Q}$. Единица общая — число $1$.
2. \textit{Отличаются:} В кольце $M_2(\mathbb{R})$ (единица $E = \text{diag}(1,1)$) возьмем подкольцо $S$ матриц вида $\begin{pmatrix} a & 0 \\ 0 & 0 \end{pmatrix}$. Внутри $S$ матрица $E' = \begin{pmatrix} 1 & 0 \\ 0 & 0 \end{pmatrix}$ работает как единица ($A E' = E' A = A, \forall A \in S$), но $E' \neq E$.

%──────────────────────────────────────────────────────────────────────────────
\section{Гомоморфизм колец. Четыре свойства. Ядро — идеал}
\addcontentsline{toc}{section}{Вопрос 3. Гомоморфизм колец. Четыре свойства. Ядро — идеал}

\begin{definition}
	Отображение $\phi: R \to S$ называется \textbf{гомоморфизмом колец}, если для любых $a, b \in R$:
	$\phi(a+b) = \phi(a) + \phi(b)$ и $\phi(ab) = \phi(a)\phi(b)$.
\end{definition}

\begin{theorem}[Четыре свойства гомоморфизма]
	Пусть $\phi: R \to S$ — гомоморфизм колец. Тогда:
	1) $\phi(0_R) = 0_S$.
	2) $\phi(-a) = -\phi(a)$.
	3) $\phi(a-b) = \phi(a) - \phi(b)$.
	4) Если $R'$ — подкольцо в $R$, то $\phi(R')$ — подкольцо в $S$.
\end{theorem}
\begin{proof}
	1) $\phi(0) = \phi(0+0) = \phi(0) + \phi(0)$, сокращая на $\phi(0)$ в группе $(S, +)$, получаем $\phi(0) = 0$.
	2) $\phi(a) + \phi(-a) = \phi(a + (-a)) = \phi(0) = 0 \Rightarrow \phi(-a) = -\phi(a)$.
	3) Стыкуется из пунктов 1 и 2.
	4) Проверяется через замкнутость операций на образе отображения.
\end{proof}

\begin{theorem}
	Ядро гомоморфизма $\Ker\phi = \{a \in R \mid \phi(a) = 0\}$ является двусторонним идеалом кольца $R$.
\end{theorem}
\begin{proof}
	Пусть $x, y \in \Ker\phi$ и $r \in R$.
	1) $\phi(x-y) = \phi(x) - \phi(y) = 0 - 0 = 0 \Rightarrow x-y \in \Ker\phi$.
	2) $\phi(rx) = \phi(r)\phi(x) = \phi(r)\cdot 0 = 0 \Rightarrow rx \in \Ker\phi$.
	3) $\phi(xr) = \phi(x)\phi(r) = 0 \cdot \phi(r) = 0 \Rightarrow xr \in \Ker\phi$. Идеал двусторонний.
\end{proof}

%──────────────────────────────────────────────────────────────────────────────
\section{Делимость в кольцах. Ассоциированные и неразложимые элементы. ООР}
\addcontentsline{toc}{section}{Вопрос 4. Делимость в кольцах. Ассоциированные и неразложимые элементы. ООР}

\begin{definition}
	Пусть $R$ --- область целостности. Говорят, что элемент $a \in R$
	\textbf{делит} элемент $b \in R$ (обозначение $a \mid b$), если существует
	$c \in R$ такое, что
	\[
	b=ac.
	\]
\end{definition}

\begin{definition}
	Элемент $u \in R$ называется \textbf{обратимым}, если существует элемент
	$v \in R$ такой, что
	\[
	uv=vu=1.
	\]
	Элемент $v$ называется обратным к $u$ и обозначается $u^{-1}$.
\end{definition}

\begin{definition}
	В области целостности элементы $a$ и $b$ называются
	\textbf{ассоциированными} ($a \sim b$), если
	\[
	a=bu,
	\]
	где $u$ --- обратимый элемент.
\end{definition}

\begin{proposition}
	Отношение ассоциированности является отношением эквивалентности.
\end{proposition}

\begin{proof}
	\begin{enumerate}
		\item Рефлексивность:
		\[
		a=a\cdot 1,
		\]
		а $1$ обратим, значит $a \sim a$.
		
		\item Симметричность: если
		\[
		a=bu,
		\]
		где $u$ обратим, то
		\[
		b=au^{-1},
		\]
		следовательно, $b \sim a$.
		
		\item Транзитивность: если
		\[
		a=bu,\qquad b=cv,
		\]
		где $u,v$ обратимы, то
		\[
		a=c(vu).
		\]
		Произведение обратимых элементов обратимо, значит $a \sim c$.
	\end{enumerate}
\end{proof}

\begin{definition}
	Делитель элемента $a$ называется \textbf{несобственным}, если он
	ассоциирован с $a$ или является обратимым элементом.
\end{definition}

\begin{definition}
	Ненулевой необратимый элемент $p$ называется \textbf{неразложимым}, если
	из равенства
	\[
	p=ab
	\]
	следует, что либо $a$, либо $b$ обратим.
\end{definition}

\begin{proposition}
	Элемент $p$ неразложим тогда и только тогда, когда у него нет собственных
	делителей.
\end{proposition}

\begin{proof}
	Пусть $p$ неразложим и
	\[
	p=ab.
	\]
	Тогда либо $a$, либо $b$ обратим.
	
	Если $a$ обратим, то
	\[
	b=a^{-1}p,
	\]
	следовательно, $b \sim p$.
	
	Аналогично, если $b$ обратим, то $a \sim p$.
	Значит, все делители элемента $p$ несобственные.
	
	Обратно, пусть у $p$ нет собственных делителей и
	\[
	p=ab.
	\]
	Тогда $a \mid p$, значит либо $a$ обратим, либо $a \sim p$.
	
	Если $a \sim p$, то
	\[
	a=pu
	\]
	для некоторого обратимого $u$. Тогда
	\[
	p=ab=pub.
	\]
	Так как область целостности не имеет делителей нуля и $p \neq 0$, получаем
	\[
	1=ub.
	\]
	Следовательно, $b$ обратим.
	
	Значит, в любом разложении $p=ab$ один из множителей обратим, то есть
	$p$ неразложим.
\end{proof}

\begin{definition}
	Область целостности называется \textbf{областью с однозначным разложением}
	(ООР, факториальным кольцом), если любой ненулевой необратимый элемент
	раскладывается в произведение неразложимых, и это разложение единственно
	с точностью до порядка множителей и ассоциированности.
\end{definition}

\begin{example}
	Кольцо целых чисел $\mathbb{Z}$ является ООР.
	Например,
	\[
	30=2 \cdot 3 \cdot 5.
	\]
	Любое другое разложение числа $30$ на неразложимые множители отличается
	только порядком и умножением на обратимые элементы $\pm 1$.
\end{example}

\begin{example}
	Кольцо многочленов $K[x]$ над полем $K$ является ООР.
\end{example}

\begin{theorem}
	Для области целостности $R$ следующие определения ООР эквивалентны:
	\begin{enumerate}
		\item Каждый ненулевой необратимый элемент раскладывается в произведение
		неразложимых, и разложение единственно с точностью до порядка и
		ассоциированности.
		
		\item Каждый ненулевой необратимый элемент раскладывается в произведение
		неразложимых, и всякий неразложимый элемент является простым.
	\end{enumerate}
\end{theorem}

\begin{proof}
	Докажем $(1)\Rightarrow(2)$.
	
	Пусть $p$ --- неразложимый элемент и
	\[
	p \mid ab.
	\]
	Тогда
	\[
	ab=pc
	\]
	для некоторого $c \in R$.
	
	Разложим элементы $a,b,c$ и $p$ на неразложимые множители.
	Так как разложение единственно, множитель $p$ должен быть ассоциирован
	с одним из неразложимых множителей произведения $ab$.
	Следовательно, $p$ ассоциирован либо с множителем из разложения $a$,
	либо с множителем из разложения $b$.
	
	Тогда
	\[
	p \mid a
	\quad \text{или} \quad
	p \mid b.
	\]
	Значит, $p$ прост.
	
	Докажем $(2)\Rightarrow(1)$.
	
	Существование разложения уже предполагается. Докажем единственность.
	
	Пусть
	\[
	p_1\cdots p_n=q_1\cdots q_m,
	\]
	где все $p_i$ и $q_j$ неразложимы.
	
	Так как $p_1$ прост и
	\[
	p_1 \mid q_1\cdots q_m,
	\]
	то
	\[
	p_1 \mid q_j
	\]
	для некоторого $j$.
	
	Но $q_j$ неразложим, следовательно, $p_1 \sim q_j$.
	
	Переставив множители, можно считать, что $j=1$.
	Тогда
	\[
	q_1=p_1u
	\]
	для некоторого обратимого $u$.
	
	Подставляя, получаем
	\[
	p_1\cdots p_n=p_1u q_2\cdots q_m.
	\]
	Сокращая на $p_1$, имеем
	\[
	p_2\cdots p_n=u q_2\cdots q_m.
	\]
	Так как $u$ обратим, его можно поглотить одним из множителей.
	
	Продолжая процесс, получаем:
	\[
	n=m,
	\]
	и после перестановки множителей
	\[
	p_i \sim q_i.
	\]
	
	Следовательно, разложение единственно с точностью до порядка и
	ассоциированности.
\end{proof}

%──────────────────────────────────────────────────────────────────────────────
\section{Взаимно простые элементы. Евклидова область. Следствия}
\addcontentsline{toc}{section}{Вопрос 5. Взаимно простые элементы. Евклидова область. Следствия}

\begin{definition}
	Элементы $a,b \in R$ называются \textbf{взаимно простыми}, если их наибольший общий делитель обратим:
	\[
	(a,b)=1.
	\]
\end{definition}

\begin{theorem}[Соотношение Безу]
	Для элементов $a,b \in R$ следующие условия эквивалентны:
	\begin{enumerate}
		\item элементы $a$ и $b$ взаимно просты;
		\item существуют $x,y \in R$ такие, что
		\[
		xa+yb=1.
		\]
	\end{enumerate}
\end{theorem}

\begin{proof}
	Пусть
	\[
	xa+yb=1.
	\]
	Если $d$ --- общий делитель элементов $a$ и $b$, то
	\[
	d \mid xa,\qquad d \mid yb,
	\]
	следовательно,
	\[
	d \mid (xa+yb)=1.
	\]
	Значит, $d$ обратим, то есть $(a,b)=1$.
	
	Обратно, пусть $(a,b)=1$.
	Тогда по соотношению Безу для НОД существуют $x,y \in R$ такие, что
	\[
	xa+yb=(a,b)=1.
	\]
\end{proof}

\begin{theorem}
	Если элементы $a$ и $b$ взаимно просты и
	\[
	a \mid bc,
	\]
	то
	\[
	a \mid c.
	\]
\end{theorem}

\begin{proof}
	Так как $(a,b)=1$, существуют $x,y \in R$ такие, что
	\[
	xa+yb=1.
	\]
	Умножим это равенство на $c$:
	\[
	xac+ybc=c.
	\]
	
	Так как
	\[
	a \mid xac
	\]
	и по условию
	\[
	a \mid bc,
	\]
	то
	\[
	a \mid ybc.
	\]
	Следовательно,
	\[
	a \mid (xac+ybc)=c.
	\]
\end{proof}

\begin{corollary}
	Если элемент $a$ взаимно прост с каждым из элементов
	\[
	b_1,b_2,\ldots,b_n,
	\]
	то
	\[
	a
	\]
	взаимно прост с произведением
	\[
	b_1b_2\cdots b_n.
	\]
\end{corollary}

\begin{proof}
	Докажем индукцией по $n$.
	
	При $n=2$ утверждение следует из предыдущей теоремы.
	
	Пусть утверждение верно для $n-1$ элементов.
	Тогда
	\[
	(a,b_1b_2\cdots b_{n-1})=1.
	\]
	Так как также
	\[
	(a,b_n)=1,
	\]
	то
	\[
	(a,b_1b_2\cdots b_n)=1.
	\]
\end{proof}

\begin{definition}
	Область целостности $R$ называется \textbf{евклидовой}, если существует функция
	\[
	N: R\setminus\{0\}\to\mathbb{Z}_{\ge0}
	\]
	(евклидова норма), такая что:
	\begin{enumerate}
		\item для любых ненулевых $a,b \in R$
		\[
		N(a)\le N(ab);
		\]
		
		\item для любых $a,b \in R$, $b\ne0$, существуют элементы
		$q,r \in R$ такие, что
		\[
		a=bq+r,
		\]
		где либо $r=0$, либо
		\[
		N(r)<N(b).
		\]
	\end{enumerate}
\end{definition}

\begin{example}
	Кольцо целых чисел $\mathbb{Z}$ является евклидовой областью
	с нормой
	\[
	N(a)=|a|.
	\]
\end{example}

\begin{example}
	Кольцо многочленов $K[x]$ над полем $K$ является евклидовой областью
	с нормой
	\[
	N(f)=\deg f.
	\]
\end{example}

\begin{theorem}[Следствия из определения евклидовой области]
	Для любой евклидовой области справедливы следующие утверждения:
	\begin{enumerate}
		\item существует алгоритм Евклида для нахождения НОД;
		\item любой идеал является главным;
		\item элемент имеет минимальную норму тогда и только тогда,
		когда он обратим;
		\item любые два элемента имеют НОД, представимый в виде
		\[
		d=xa+yb.
		\]
	\end{enumerate}
\end{theorem}

\begin{proof}
	\begin{enumerate}
		\item Пусть заданы элементы $a,b$ ($b\ne0$).
		По определению евклидовой области:
		\[
		a=bq_1+r_1,
		\qquad N(r_1)<N(b).
		\]
		Далее:
		\[
		b=r_1q_2+r_2,
		\qquad N(r_2)<N(r_1),
		\]
		и т.д.
		
		Получаем последовательность норм:
		\[
		N(b)>N(r_1)>N(r_2)>\cdots
		\]
		Она строго убывает и состоит из неотрицательных целых чисел,
		поэтому процесс конечен.
		
		Последний ненулевой остаток и есть НОД элементов $a$ и $b$.
		
		\item Пусть $I$ --- ненулевой идеал.
		Выберем элемент $d\in I$ минимальной нормы.
		
		Докажем, что
		\[
		I=(d).
		\]
		
		Пусть $a\in I$.
		Разделим $a$ на $d$:
		\[
		a=dq+r,
		\qquad r=0
		\quad \text{или} \quad
		N(r)<N(d).
		\]
		
		Так как
		\[
		r=a-dq,
		\]
		а идеал замкнут относительно сложения и умножения,
		то
		\[
		r\in I.
		\]
		
		Минимальность нормы элемента $d$ означает, что случай
		\[
		N(r)<N(d)
		\]
		невозможен для ненулевого $r$.
		Следовательно,
		\[
		r=0.
		\]
		
		Значит,
		\[
		a=dq,
		\]
		то есть $a\in(d)$.
		
		Следовательно,
		\[
		I=(d).
		\]
		
		\item Пусть $u$ обратим.
		Тогда существует элемент $v$ такой, что
		\[
		uv=1.
		\]
		
		По свойству нормы:
		\[
		N(u)\le N(uv)=N(1).
		\]
		Следовательно, обратимые элементы имеют минимальную норму.
		
		Обратно, пусть $a$ имеет минимальную норму.
		Разделим $1$ на $a$:
		\[
		1=aq+r,
		\qquad r=0
		\quad \text{или} \quad
		N(r)<N(a).
		\]
		
		Минимальность нормы исключает возможность ненулевого остатка.
		Поэтому
		\[
		r=0,
		\]
		и
		\[
		1=aq.
		\]
		
		Следовательно, $a$ обратим.
		
		\item Рассмотрим множество
		\[
		I=\{xa+yb \mid x,y\in R\}.
		\]
		Это идеал.
		
		По пункту 2:
		\[
		I=(d)
		\]
		для некоторого элемента $d$.
		
		Так как
		\[
		a,b\in I,
		\]
		то
		\[
		d\mid a,
		\qquad d\mid b.
		\]
		
		Если $c$ --- любой общий делитель элементов $a$ и $b$, то
		\[
		c\mid xa+yb
		\]
		для любых $x,y$.
		Следовательно,
		\[
		c\mid d.
		\]
		
		Значит, $d$ является НОД элементов $a$ и $b$.
		
		По определению идеала:
		\[
		d=xa+yb.
		\]
	\end{enumerate}
\end{proof}

%──────────────────────────────────────────────────────────────────────────────
\section{Евклидова область является ООР}
\addcontentsline{toc}{section}{Вопрос 6. Евклидова область является ООР}

\begin{theorem}
	Всякая евклидова область является областью с однозначным разложением (ООР).
\end{theorem}
\begin{proof}
	Полноценное доказательство состоит из двух шагов:
	1) \textbf{Существование разложения:} Если элемент неразложим, разложение есть. Если разложим, делим на множители. Так как при делении норма строго уменьшается, цепочка делителей не может быть бесконечной, процесс оборвется на неразложимых элементах.
	2) \textbf{Единственность:} Базируется на свойстве, что если неразложимый элемент $p$ делит произведение $ab$, то он делит $a$ или $b$ (что доказывается через линейное представление НОД в евклидовых кольцах). Далее применяется стандартная индукция по числу множителей.
\end{proof}

%──────────────────────────────────────────────────────────────────────────────
\section{Кольцо полиномов $R[x]$}
\addcontentsline{toc}{section}{Вопрос 7. Кольцо полиномов $R[x]$}

\begin{definition}
	Множество финитных последовательностей элементов кольца $R$
	вида
	\[
	(a_0,a_1,\dots,a_n,0,0,\dots)
	\]
	с покомпонентным сложением и умножением по правилу Коши
	\[
	c_k=\sum_{i+j=k} a_i b_j
	\]
	называется \textbf{кольцом полиномов} и обозначается $R[x]$.
\end{definition}

\begin{definition}
	Элемент
	\[
	f=(a_0,a_1,\dots,a_n,0,\dots)
	\]
	записывают в виде
	\[
	f(x)=a_0+a_1x+\dots+a_nx^n.
	\]
	Числа $a_i$ называются коэффициентами полинома.
\end{definition}

\begin{definition}
	Одночлен
	\[
	x^n
	\]
	представляется последовательностью
	\[
	(0,\dots,0,1,0,\dots),
	\]
	где единица стоит на $n$-й позиции.
\end{definition}

\begin{example}
	\[
	1=(1,0,0,\dots),
	\]
	\[
	x=(0,1,0,\dots),
	\]
	\[
	x^2=(0,0,1,0,\dots).
	\]
\end{example}

\begin{theorem}
	Множество $R[x]$ является кольцом.
\end{theorem}

\begin{proof}
	Проверим аксиомы кольца.
	
	\begin{enumerate}
		\item Замкнутость относительно сложения.
		
		Если
		\[
		f=(a_0,a_1,\dots),
		\qquad
		g=(b_0,b_1,\dots),
		\]
		то
		\[
		f+g=(a_0+b_0,a_1+b_1,\dots).
		\]
		Так как последовательности финитны, сумма также финитна.
		
		\item Ассоциативность сложения следует из ассоциативности сложения в $R$.
		
		\item Нулевой элемент:
		\[
		0=(0,0,\dots).
		\]
		
		\item Противоположный элемент:
		\[
		-f=(-a_0,-a_1,\dots).
		\]
		
		\item Коммутативность сложения следует из коммутативности сложения в $R$.
		
		\item Замкнутость относительно умножения.
		
		Для
		\[
		f=(a_i),\qquad g=(b_i)
		\]
		определим произведение:
		\[
		fg=(c_k),
		\qquad
		c_k=\sum_{i+j=k} a_i b_j.
		\]
		
		Так как последовательности финитны, в сумме лишь конечное число
		ненулевых слагаемых, поэтому произведение снова финитно.
		
		\item Ассоциативность умножения следует из ассоциативности умножения в $R$:
		\[
		((fg)h)_n
		=
		\sum_{i+j+k=n} a_i b_j c_k
		=
		(f(gh))_n.
		\]
		
		\item Дистрибутивность:
		\[
		f(g+h)=fg+fh,
		\]
		что проверяется раскрытием сумм.
		
	\end{enumerate}
	
	Следовательно, $R[x]$ является кольцом.
\end{proof}

\begin{theorem}
	Если кольцо $R$ имеет единицу, то $R[x]$ также имеет единицу.
\end{theorem}

\begin{proof}
	Рассмотрим последовательность
	\[
	(1,0,0,\dots).
	\]
	
	Для любого полинома
	\[
	f=(a_0,a_1,\dots)
	\]
	имеем:
	\[
	(1,0,\dots)\cdot f=f.
	\]
	
	Следовательно,
	\[
	1_{R[x]}=(1,0,0,\dots).
	\]
\end{proof}

\begin{theorem}
	Кольцо $R$ вкладывается в кольцо $R[x]$.
\end{theorem}

\begin{proof}
	Определим отображение
	\[
	\varphi:R\to R[x]
	\]
	по правилу
	\[
	\varphi(a)=(a,0,0,\dots).
	\]
	
	Проверим сохранение операций:
	\[
	\varphi(a+b)
	=
	(a+b,0,\dots)
	=
	(a,0,\dots)+(b,0,\dots)
	=
	\varphi(a)+\varphi(b),
	\]
	\[
	\varphi(ab)
	=
	(ab,0,\dots)
	=
	(a,0,\dots)(b,0,\dots)
	=
	\varphi(a)\varphi(b).
	\]
	
	Если
	\[
	\varphi(a)=\varphi(b),
	\]
	то
	\[
	(a,0,\dots)=(b,0,\dots),
	\]
	следовательно,
	\[
	a=b.
	\]
	
	Значит, отображение инъективно, то есть $R$ изоморфно
	подкольцу константных полиномов в $R[x]$.
\end{proof}

\begin{theorem}
	Если $R$ --- область целостности, то $R[x]$ также является областью целостности.
\end{theorem}

\begin{proof}
	Пусть
	\[
	f(x)=a_nx^n+\dots+a_0,
	\qquad
	g(x)=b_mx^m+\dots+b_0
	\]
	--- ненулевые полиномы.
	
	Тогда
	\[
	a_n\ne0,
	\qquad
	b_m\ne0.
	\]
	
	Рассмотрим коэффициент при старшей степени в произведении:
	\[
	f(x)g(x).
	\]
	
	Коэффициент при
	\[
	x^{n+m}
	\]
	равен
	\[
	a_n b_m.
	\]
	
	Так как $R$ --- область целостности, делителей нуля нет, поэтому
	\[
	a_n b_m\ne0.
	\]
	
	Следовательно,
	\[
	f(x)g(x)\ne0.
	\]
	
	Значит, в $R[x]$ нет делителей нуля, то есть $R[x]$
	является областью целостности.
\end{proof}

%──────────────────────────────────────────────────────────────────────────────
\section{Подстановка в полином. Алгебраические и трансцендентные элементы}
\addcontentsline{toc}{section}{Вопрос 8. Подстановка в полином. Алгебраические и трансцендентные элементы}

Пусть $R$ --- кольцо, $S$ --- унитарное надкольцо кольца $R$,
$x \in S$.

Для полинома
\[
f(t)=a_0+a_1t+\dots+a_nt^n \in R[t]
\]
определим значение полинома в точке $x$:
\[
f(x)=a_0+a_1x+\dots+a_nx^n.
\]

\begin{theorem}
	Отображение
	\[
	\phi_x:R[t]\to S,
	\qquad
	\phi_x(f(t))=f(x),
	\]
	является гомоморфизмом колец.
\end{theorem}

\begin{proof}
	Пусть
	\[
	f(t)=\sum_{i=0}^n a_it^i,
	\qquad
	g(t)=\sum_{i=0}^m b_it^i.
	\]
	
	Проверим сохранение сложения:
	\[
	\phi_x(f+g)
	=
	(f+g)(x)
	=
	\sum (a_i+b_i)x^i.
	\]
	
	Тогда
	\[
	(f+g)(x)
	=
	\sum a_ix^i+\sum b_ix^i
	=
	f(x)+g(x).
	\]
	
	Следовательно,
	\[
	\phi_x(f+g)=\phi_x(f)+\phi_x(g).
	\]
	
	Проверим сохранение умножения:
	\[
	f(t)g(t)=\sum_k c_k t^k,
	\qquad
	c_k=\sum_{i+j=k} a_i b_j.
	\]
	
	Тогда
	\[
	\phi_x(fg)
	=
	(fg)(x)
	=
	\sum_k c_k x^k.
	\]
	
	Подставляя выражение для $c_k$, получаем
	\[
	(fg)(x)
	=
	\sum_k \left(\sum_{i+j=k} a_i b_j\right)x^k.
	\]
	
	Так как
	\[
	x^k=x^{i+j}=x^i x^j,
	\]
	имеем
	\[
	(fg)(x)
	=
	\left(\sum_i a_i x^i\right)
	\left(\sum_j b_j x^j\right).
	\]
	
	Следовательно,
	\[
	\phi_x(fg)=\phi_x(f)\phi_x(g).
	\]
	
	Кроме того,
	\[
	\phi_x(1)=1.
	\]
	
	Следовательно, $\phi_x$ является гомоморфизмом колец.
\end{proof}

\begin{definition}
	Элемент $x \in S$ называется \textbf{алгебраическим} над $R$,
	если существует ненулевой многочлен
	\[
	f(t)\in R[t]
	\]
	такой, что
	\[
	f(x)=0.
	\]
	
	Эквивалентно:
	\[
	\Ker \phi_x \ne \{0\}.
	\]
\end{definition}

\begin{definition}
	Если
	\[
	\Ker \phi_x=\{0\},
	\]
	то элемент $x$ называется \textbf{трансцендентным} над $R$.
\end{definition}

\begin{example}
	Число
	\[
	\sqrt{2}
	\]
	алгебраично над $\mathbb{Q}$, так как
	\[
	x^2-2=0.
	\]
\end{example}

\begin{example}
	Число
	\[
	\pi
	\]
	трансцендентно над $\mathbb{Q}$.
\end{example}

\begin{theorem}
	Все кольца полиномов над кольцом $R$ изоморфны.
\end{theorem}

\begin{proof}
	Пусть $R[x]$ и $R[y]$ --- два кольца полиномов над $R$
	с различными обозначениями переменной.
	
	Определим отображение
	\[
	\varphi:R[x]\to R[y]
	\]
	по правилу
	\[
	\varphi\left(
	\sum_{i=0}^n a_i x^i
	\right)
	=
	\sum_{i=0}^n a_i y^i.
	\]
	
	Проверим сохранение сложения:
	\[
	\varphi(f+g)=\varphi(f)+\varphi(g).
	\]
	
	Действительно, коэффициенты складываются одинаково.
	
	Проверим сохранение умножения:
	\[
	\varphi(fg)=\varphi(f)\varphi(g),
	\]
	так как правило умножения коэффициентов в $R[x]$
	и $R[y]$ совпадает.
	
	Отображение биективно, поскольку каждому полиному
	\[
	\sum a_i y^i
	\]
	соответствует единственный полином
	\[
	\sum a_i x^i.
	\]
	
	Следовательно,
	\[
	R[x]\cong R[y].
	\]
\end{proof}

%──────────────────────────────────────────────────────────────────────────────
\section{Конечная область целостности — поле. Вложение кольца в поле}
\addcontentsline{toc}{section}{Вопрос 9. Конечная область целостности — поле. Вложение кольца в поле}

\begin{theorem}
	Всякая конечная область целостности является полем.
\end{theorem}

\begin{proof}
	Пусть
	\[
	R=\{x_1,x_2,\dots,x_n\}
	\]
	--- конечная область целостности.
	
	Возьмем произвольный ненулевой элемент
	\[
	a\in R.
	\]
	
	Рассмотрим отображение
	\[
	f:R\to R,
	\qquad
	f(x)=ax.
	\]
	
	Докажем, что оно инъективно.
	
	Пусть
	\[
	ax_i=ax_j.
	\]
	Тогда
	\[
	a(x_i-x_j)=0.
	\]
	
	Так как $R$ --- область целостности и
	\[
	a\ne0,
	\]
	то делителей нуля нет, следовательно,
	\[
	x_i-x_j=0,
	\]
	то есть
	\[
	x_i=x_j.
	\]
	
	Значит, отображение $f$ инъективно.
	
	Так как множество $R$ конечно, инъективное отображение
	\[
	f:R\to R
	\]
	является сюръективным.
	
	Следовательно, существует элемент
	\[
	b\in R
	\]
	такой, что
	\[
	ab=1.
	\]
	
	Значит, элемент $a$ обратим.
	
	Так как произвольный ненулевой элемент обратим,
	$R$ является полем.
\end{proof}

\begin{theorem}
	Если кольцо $R$ вкладывается в поле, то $R$ является областью целостности.
\end{theorem}

\begin{proof}
	Пусть существует инъективный гомоморфизм
	\[
	\varphi:R\to K,
	\]
	где $K$ --- поле.
	
	Предположим, что
	\[
	ab=0
	\]
	для некоторых
	\[
	a,b\in R.
	\]
	
	Тогда
	\[
	\varphi(ab)=\varphi(0)=0.
	\]
	
	Так как $\varphi$ --- гомоморфизм,
	\[
	\varphi(a)\varphi(b)=0.
	\]
	
	Поле не имеет делителей нуля, поэтому
	\[
	\varphi(a)=0
	\quad \text{или} \quad
	\varphi(b)=0.
	\]
	
	Так как $\varphi$ инъективно,
	\[
	a=0
	\quad \text{или} \quad
	b=0.
	\]
	
	Следовательно, в $R$ нет делителей нуля, то есть
	$R$ является областью целостности.
\end{proof}

\begin{theorem}
	Всякая область целостности вкладывается в некоторое поле.
\end{theorem}

\begin{proof}
	Пусть $R$ --- область целостности.
	
	Рассмотрим множество пар
	\[
	(a,b),
	\qquad
	a\in R,\ b\in R\setminus\{0\}.
	\]
	
	Введем отношение эквивалентности:
	\[
	(a,b)\sim(c,d)
	\iff
	ad=bc.
	\]
	
	Обозначим класс эквивалентности пары $(a,b)$ через
	\[
	\frac{a}{b}.
	\]
	
	Множество всех таких классов обозначим через $K$.
	
	Определим операции:
	\[
	\frac{a}{b}+\frac{c}{d}
	=
	\frac{ad+bc}{bd},
	\]
	\[
	\frac{a}{b}\cdot\frac{c}{d}
	=
	\frac{ac}{bd}.
	\]
	
	Корректность операций проверяется непосредственно.
	
	Множество $K$ является полем.
	Действительно, для любого ненулевого элемента
	\[
	\frac{a}{b},
	\qquad a\ne0,
	\]
	обратным является элемент
	\[
	\frac{b}{a},
	\]
	так как
	\[
	\frac{a}{b}\cdot\frac{b}{a}
	=
	\frac{ab}{ba}
	=
	1.
	\]
	
	Определим отображение
	\[
	\varphi:R\to K
	\]
	по правилу
	\[
	\varphi(a)=\frac{a}{1}.
	\]
	
	Проверим инъективность.
	
	Если
	\[
	\frac{a}{1}=\frac{b}{1},
	\]
	то
	\[
	a\cdot1=b\cdot1,
	\]
	следовательно,
	\[
	a=b.
	\]
	
	Значит, $\varphi$ --- инъективный гомоморфизм.
	
	Следовательно, область целостности $R$
	вкладывается в поле $K$.
\end{proof}

\begin{definition}
	Поле $K$, построенное из области целостности $R$ указанным способом,
	называется \textbf{полем частных} области целостности $R$.
\end{definition}
%──────────────────────────────────────────────────────────────────────────────
\section{Поле частных области целостности}
\addcontentsline{toc}{section}{Вопрос 10. Поле частных области целостности}

\begin{theorem}
	Всякая область целостности может быть вложена в поле.
\end{theorem}

\begin{proof}
	Пусть $R$ --- область целостности.
	
	Рассмотрим множество пар
	\[
	(a,b),
	\qquad
	a\in R,\quad b\in R\setminus\{0\}.
	\]
	
	Введем отношение эквивалентности:
	\[
	(a,b)\sim(c,d)
	\iff
	ad=bc.
	\]
	
	Докажем, что это действительно отношение эквивалентности.
	
	\begin{enumerate}
		\item Рефлексивность:
		\[
		ab=ba,
		\]
		следовательно,
		\[
		(a,b)\sim(a,b).
		\]
		
		\item Симметричность:
		если
		\[
		ad=bc,
		\]
		то
		\[
		bc=ad,
		\]
		следовательно,
		\[
		(c,d)\sim(a,b).
		\]
		
		\item Транзитивность:
		
		Пусть
		\[
		(a,b)\sim(c,d),
		\qquad
		(c,d)\sim(e,f).
		\]
		
		Тогда
		\[
		ad=bc,
		\qquad
		cf=de.
		\]
		
		Умножая первое равенство на $f$, второе --- на $b$,
		получаем:
		\[
		adf=bcf,
		\]
		\[
		bcf=bde.
		\]
		
		Следовательно,
		\[
		adf=bde.
		\]
		
		Так как $R$ --- область целостности и $d\ne0$,
		можно сократить на $d$:
		\[
		af=be.
		\]
		
		Следовательно,
		\[
		(a,b)\sim(e,f).
		\]
	\end{enumerate}
\end{proof}

\begin{definition}
	Класс эквивалентности пары $(a,b)$ обозначается:
	\[
	\frac{a}{b}.
	\]
\end{definition}

\begin{definition}
	Множество всех классов эквивалентности обозначается
	\[
	\operatorname{Frac}(R)
	\]
	и называется \textbf{полем частных} области целостности $R$.
\end{definition}

\begin{theorem}
	На множестве $\operatorname{Frac}(R)$ корректно определены операции:
	\[
	\frac{a}{b}+\frac{c}{d}
	=
	\frac{ad+bc}{bd},
	\]
	\[
	\frac{a}{b}\cdot\frac{c}{d}
	=
	\frac{ac}{bd}.
	\]
\end{theorem}

\begin{proof}
	Докажем корректность сложения.
	
	Пусть
	\[
	\frac{a}{b}=\frac{a'}{b'},
	\qquad
	\frac{c}{d}=\frac{c'}{d'}.
	\]
	
	Тогда
	\[
	ab'=a'b,
	\qquad
	cd'=c'd.
	\]
	
	Нужно доказать:
	\[
	\frac{ad+bc}{bd}
	=
	\frac{a'd'+b'c'}{b'd'}.
	\]
	
	Проверим условие эквивалентности:
	\[
	(ad+bc)b'd'
	=
	adb'd'+bcb'd'.
	\]
	
	Используя
	\[
	ab'=a'b,
	\qquad
	cd'=c'd,
	\]
	получаем:
	\[
	adb'd'=a'dbd',
	\]
	\[
	bcb'd'=b'c'bd.
	\]
	
	Следовательно,
	\[
	(ad+bc)b'd'
	=
	(a'd'+b'c')bd.
	\]
	
	Значит, сложение определено корректно.
	
	Корректность умножения проверяется аналогично.
\end{proof}

\begin{theorem}
	$\operatorname{Frac}(R)$ является полем.
\end{theorem}

\begin{proof}
	Аксиомы кольца проверяются непосредственно.
	
	Нулевой элемент:
	\[
	0=\frac{0}{1}.
	\]
	
	Единица:
	\[
	1=\frac{1}{1}.
	\]
	
	Пусть
	\[
	\frac{a}{b}\ne0.
	\]
	Тогда
	\[
	a\ne0.
	\]
	
	Рассмотрим элемент
	\[
	\frac{b}{a}.
	\]
	
	Тогда
	\[
	\frac{a}{b}\cdot\frac{b}{a}
	=
	\frac{ab}{ba}
	=
	\frac{1}{1}.
	\]
	
	Следовательно, любой ненулевой элемент обратим.
	
	Значит,
	\[
	\operatorname{Frac}(R)
	\]
	является полем.
\end{proof}

\begin{theorem}
	Область целостности $R$ вкладывается в свое поле частных.
\end{theorem}

\begin{proof}
	Определим отображение
	\[
	\varphi:R\to\operatorname{Frac}(R)
	\]
	по правилу
	\[
	\varphi(a)=\frac{a}{1}.
	\]
	
	Проверим сохранение операций:
	\[
	\varphi(a+b)
	=
	\frac{a+b}{1}
	=
	\frac{a}{1}+\frac{b}{1},
	\]
	\[
	\varphi(ab)
	=
	\frac{ab}{1}
	=
	\frac{a}{1}\cdot\frac{b}{1}.
	\]
	
	Если
	\[
	\varphi(a)=\varphi(b),
	\]
	то
	\[
	\frac{a}{1}=\frac{b}{1}.
	\]
	
	Следовательно,
	\[
	a\cdot1=b\cdot1,
	\]
	то есть
	\[
	a=b.
	\]
	
	Значит, отображение инъективно.
\end{proof}

\begin{example}
	Полем частных кольца целых чисел $\mathbb{Z}$ является поле
	рациональных чисел:
	\[
	\operatorname{Frac}(\mathbb{Z})=\mathbb{Q}.
	\]
	
	Элементами являются дроби вида
	\[
	\frac{a}{b},
	\qquad
	a,b\in\mathbb{Z},\quad b\ne0.
	\]
\end{example}

\begin{example}
	Кольцо гауссовых целых чисел
	\[
	\mathbb{Z}[i]
	=
	\{a+bi\mid a,b\in\mathbb{Z}\}
	\]
	является областью целостности.
	
	Его полем частных является поле гауссовых рациональных чисел:
	\[
	\mathbb{Q}(i)
	=
	\{a+bi\mid a,b\in\mathbb{Q}\}.
	\]
\end{example}
%──────────────────────────────────────────────────────────────────────────────
\section{Поле рациональных дробей $F(x)$. Бесконечность и характеристика}
\addcontentsline{toc}{section}{Вопрос 11. Поле рациональных дробей $F(x)$. Бесконечность и характеристика}

\begin{theorem}
	Поле $F(x)$ рациональных дробей над любым полем $F$ всегда бесконечно.
\end{theorem}
\begin{proof}
	Рассмотрим бесконечную последовательность элементов вида $x, x^2, x^3, \dots, x^n, \dots$. Все эти элементы различны в $F(x)$, поскольку многочлены разной степени не могут быть равны как финитные последовательности. Следовательно, поле содержит бесконечно много элементов.
\end{proof}

Характеристика поля $F(x)$ совпадает с характеристикой исходного поля констант $F$, так как единица поля $F(x)$ совпадает с единицей поля $F$. Если $\operatorname{char} F = p$, то $p \cdot 1_{F(x)} = p \cdot 1_F = 0$.

%──────────────────────────────────────────────────────────────────────────────
\section{Делители, обратимые, ассоциированные элементы в $K[x]$. Наибольший общий делитель (НОД) и наименьшее общее кратное (НОК), их свойства}
\addcontentsline{toc}{section}{Вопрос 12. Делители, обратимые, ассоциированные элементы в $K[x]$. Наибольший общий делитель (НОД) и наименьшее общее кратное (НОК), их свойства}

Пусть $K$ — произвольное поле, а $K[x]$ — кольцо многочленов от одной переменной над этим полем. Поскольку $K$ является полем, кольцо $K[x]$ является евклидовой областью, а следовательно — главным идеалом и областью с однозначным разложением (ООР).

\begin{definition}
	Многочлен $g(x)$ называется \textbf{делителем} многочлена $f(x)$ (обозначается $g \mid f$), если существует такой многочлен $q(x) \in K[x]$, что $f(x) = g(x)q(x)$.
\end{definition}

\begin{definition}
	Элемент $f(x) \in K[x]$ называется \textbf{обратимым} (или единицей кольца), если существует $f^{-1}(x) \in K[x]$, такой что $f(x)f^{-1}(x) = 1$. Поскольку для любых многочленов $\deg(f \cdot g) = \deg f + \deg g$, обратимыми элементами в $K[x]$ являются только ненулевые константы (многочлены нулевой степени): $U(K[x]) = K^{\times} = K \setminus \{0\}$.
\end{definition}

\begin{definition}
	Многочлены $f(x)$ and $g(x)$ называются \textbf{ассоциированными} ($f \sim g$), если они делят друг друга ($f \mid g$ и $g \mid f$). В области целостности это эквивалентно тому, что они отличаются на обратимый элемент кольца: $f(x) = c \cdot g(x)$, где $c \in K^{\times}$.
\end{definition}

\begin{definition}
	Многочлен называется \textbf{унитарным} (или нормированным), если его старший коэффициент равен $1$. В каждом классе ассоциированных ненулевых многочленов существует ровно один унитарный многочлен.
\end{definition}

\begin{definition}
	\textbf{Наибольшим общим делителем (НОД)} многочленов $f(x)$ и $g(x)$ называется многочлен $d(x) \in K[x]$, удовлетворяющий двум условиям:
	\begin{enumerate}
		\item $d(x)$ является их общим делителем: $d \mid f$ и $d \mid g$;
		\item $d(x)$ делится на любой другой их общий делитель: если $h \mid f$ и $h \mid g$, то $h \mid d$.
	\end{enumerate}
	Чтобы обеспечить единственность НОД, в качестве стандартного представителя выбирают \textbf{унитарный} многочлен и обозначают $\gcd(f, g)$ или $(f, g)$. Если $\gcd(f, g) = 1$, многочлены называются \textbf{взаимно простыми}.
\end{definition}

Для наибольшего общего делителя справедливы следующие фундаментальные свойства:
\begin{enumerate}
	\item \textbf{Линейное представление (соотношение Безу):} НОД двух многочленов всегда можно представить в виде их линейной комбинации: $\gcd(f, g) = u(x)f(x) + v(x)g(x)$, где $u(x), v(x) \in K[x]$.
	\item \textbf{Связь с идеалами:} Главный идеал, порожденный НОД, совпадает с суммой идеалов, порожденных самими многочленами: $\langle f \rangle + \langle g \rangle = \langle \gcd(f, g) \rangle$.
	\item \textbf{Алгоритмичность:} НОД эффективно находится с помощью алгоритма Евклида путем последовательного деления многочленов с остатком.
\end{enumerate}

\begin{definition}
	\textbf{Наименьшим общим кратным (НОК)} многочленов $f(x)$ и $g(x)$ называется многочлен $L(x) \in K[x]$, удовлетворяющий двум условиям:
	\begin{enumerate}
		\item $L(x)$ является их общим кратным: $f \mid L$ и $g \mid L$;
		\item $L(x)$ делит любое другое их общее кратное: если $f \mid M$ и $g \mid M$, то $L \mid M$.
	\end{enumerate}
	По аналогии, для единственности НОК выбирается \textbf{унитарным} многочленом и обозначается $\operatorname{lcm}(f, g)$ или $[f, g]$.
\end{definition}

Понятия НОД и НОК тесно связаны между собой. Для любых двух ненулевых многочленов $f(x)$ и $g(x)$ произведение их НОД и НОК (с точностью до умножения на унитарный коэффициент старшего члена) равно произведению самих многочленов. Если зафиксировать их унитарные представители, то выполнено строгое равенство:
\[
\gcd(f, g) \cdot \operatorname{lcm}(f, g) = f_0(x) \cdot g_0(x),
\]
где $f_0(x)$ и $g_0(x)$ — многочлены $f(x)$ и $g(x)$, приведенные к унитарному виду (то есть разделенные на свои старшие коэффициенты).

%──────────────────────────────────────────────────────────────────────────────
\section{Евклидовы кольца: определение и три примера}

\begin{definition}
	Область целостности $R$ называется \textbf{евклидовой областью}, если существует функция $N \colon R \setminus \{0\} \to \mathbb{N}_0$ (\textbf{евклидова норма}) такая, что для любых $a \in R$, $b \in R \setminus \{0\}$ существуют $q, r \in R$:
	\[
	a = qb + r, \quad \text{где } r = 0 \text{ или } N(r) < N(b).
	\]
\end{definition}

\begin{example}[$\Z$]
	Норма $N(n) = |n|$. Для любых $a, b \in \Z$, $b \neq 0$, существует деление с остатком $a = qb + r$, $0 \leq r < |b|$.
\end{example}

\begin{example}[${F[x]}$]
	Кольцо многочленов над полем $F$ с нормой $N(f) = \deg f$. Для $f, g \in F[x]$,
	$g \neq 0$, алгоритм деления «углом» даёт $f = qg + r$, где $r = 0$ или
	$\deg r < \deg g$. Деление возможно, поскольку ведущий коэффициент $g$
	обратим в поле $F$.
\end{example}

\begin{example}[{$\Z[i]$}]
	Гауссовы целые $\Z[i] = \{a + bi \mid a, b \in \Z\}$ с нормой $N(a+bi) = a^2 + b^2$.
	Для $\alpha, \beta \in \Z[i]$, $\beta \neq 0$, рассмотрим $\alpha/\beta = x + yi \in \Q(i)$.
	Выберем $m, n \in \Z$ с $|x - m| \leq \tfrac{1}{2}$, $|y - n| \leq \tfrac{1}{2}$.
	Положим $q = m + ni$, $r = \alpha - q\beta$. Тогда:
	\[
	N(r) = N(\beta) \cdot N\!\left(\tfrac{\alpha}{\beta} - q\right)
	\leq N(\beta)\!\left(\tfrac{1}{4} + \tfrac{1}{4}\right)
	= \tfrac{N(\beta)}{2} < N(\beta).
	\]
\end{example}

%──────────────────────────────────────────────────────────────────────────────
\section{Взаимно простые элементы, алгоритм Евклида и представление НОД}

\textbf{Определение.} Элементы $a, b$ кольца $R$ называются взаимно простыми, если их наибольший общий делитель $\text{НОД}(a, b) = 1$ (или любому обратимому элементу).

\textbf{Алгоритм Евклида.} Пусть $R$ — евклидово кольцо с нормой $N$. Для любых $a, b \in R$ ($b \neq 0$) можно записать цепочку делений с остатком:
\begin{align*}
	a &= q_1 b + r_1, \quad N(r_1) < N(b) \\
	b &= q_2 r_1 + r_2, \quad N(r_2) < N(r_1) \\
	&\dots \\
	r_{k-2} &= q_k r_{k-1} + r_k, \quad N(r_k) < N(r_{k-1}) \\
	r_{k-1} &= q_{k+1} r_k + 0
\end{align*}
\textbf{Доказательство конечности:} Так как $N(b) > N(r_1) > N(r_2) > \dots \geq 0$ — строго убывающая последовательность целых неотрицательных чисел, она обязана оборваться за конечное число шагов. Последний ненулевой остаток $r_k$ является $\text{НОД}(a, b)$.

\textbf{Линейное представление НОД (Тождество Безу).} 
\textbf{Теорема:} В евклидовом кольце для любых $a, b$ существуют $x, y \in R$, такие что $\text{НОД}(a, b) = xa + yb$.
\textbf{Доказательство:} Проводится обратным ходом алгоритма Евклида. Из предпоследнего уравнения выражаем $r_k = r_{k-2} - q_k r_{k-1}$. Затем подставляем $r_{k-1}$ из предыдущего уравнения и так далее, поднимаясь вверх, пока не выразим $r_k$ через исходные $a$ и $b$.

\textbf{Свойства взаимно простых элементов:}
1. Если $a$ и $b$ взаимно просты, то существуют $x, y \in R$, такие что $xa + yb = 1$.
2. Если $a \vdots c$ и $b \vdots c$, и при этом $a$ и $b$ взаимно просты, то $c$ — обратимый элемент.
3. Если $ab \vdots c$ и $\text{НОД}(a, c) = 1$, то $b \vdots c$.

%──────────────────────────────────────────────────────────────────────────────
\section{Неприводимые элементы. Теорема о факториальности евклидовых колец}

\begin{definition}
	Ненулевой необратимый элемент $p \in R$ называется \textbf{неприводимым} (irreducible), если из $p = ab$ следует $a \in R^*$ или $b \in R^*$.
\end{definition}

\begin{example}
	В $\Z$: неприводимые --- это $\pm 2, \pm 3, \pm 5, \ldots$ В $F[x]$: неприводимые многочлены. В $\Z[i]$: элементы $1+i$, $3$, $2+i$.
\end{example}

\begin{theorem}
	В ОГИ (и в евклидовой области) если $p$ неприводимый и $p \mid a_1 a_2 \cdots a_n$, то $p \mid a_i$ для некоторого $i$.
\end{theorem}
\begin{proof}
	В ОГИ идеал $(p)$ максимален (вопрос 40), а значит прост. Из $p \mid ab$ следует $ab \in (p)$, откуда $a \in (p)$ или $b \in (p)$, т.\,е.\ $p \mid a$ или $p \mid b$. Для $n$ множителей --- индукция. \qed
\end{proof}

\begin{theorem}
	Каждая евклидова область факториальна.
\end{theorem}
\begin{proof}
	Евклидова $\Rightarrow$ ОГИ (вопрос 39) $\Rightarrow$ ООР (вопрос 38). \qed
\end{proof}

%──────────────────────────────────────────────────────────────────────────────
\section{Идеал. Примеры. Радикал коммутативного кольца}
\textbf{Определение.} Непустое подмножество $I$ кольца $R$ называется идеалом (двусторонним), если:
1. $\forall a, b \in I \implies a - b \in I$ ($I$ — подгруппа по сложению).
2. $\forall a \in I, \forall r \in R \implies ra \in I$ и $ar \in I$.

\textbf{Примеры:} В кольце $\mathbb{Z}$ множество чётных чисел $2\mathbb{Z}$ является идеалом. В кольце многочленов $\mathbb{R}[x]$ множество многочленов без свободного члена $x\mathbb{R}[x]$ — идеал.

\textbf{Идеал с обратимым элементом.} Если идеал $I$ содержит обратимый элемент $u$ (т.е. $\exists u^{-1} \in R$), то $\forall r \in R$ выполнено $r = r \cdot 1 = r \cdot (u^{-1} u) = (r u^{-1}) u \in I$, так как $u \in I$. Значит, $I = R$. В частности, идеал, содержащий единицу, совпадает со всем кольцом.

\textbf{Радикал коммутативного кольца (Нильрадикал).} $\text{Nil}(R) = \{a \in R \mid \exists n \in \mathbb{N} : a^n = 0\}$ — множество всех нильпотентных элементов.

\textbf{Доказательство, что $\text{Nil}(R)$ — идеал:}
Пусть $a, b \in \text{Nil}(R)$, значит $\exists n, m \in \mathbb{N}$ такие, что $a^n = 0, b^m = 0$. 
Покажем, что $a-b \in \text{Nil}(R)$. Так как кольцо $R$ \textbf{коммутативно}, применим формулу бинома Ньютона для степени $N = n + m - 1$:
$$ (a - b)^N = \sum_{k=0}^{N} (-1)^{N-k} \binom{N}{k} a^k b^{N-k} $$
В каждом слагаемом либо $k \geq n$, и тогда $a^k = a^n a^{k-n} = 0$, либо $k < n$, но тогда $N - k = n + m - 1 - k > m - 1$, то есть $N - k \geq m$, откуда $b^{N-k} = 0$. Таким образом, все слагаемые равны нулю, и $(a-b)^N = 0 \implies a-b \in \text{Nil}(R)$.
Для умножения: если $a^n = 0$ и $r \in R$, то в силу коммутативности $(ra)^n = r^n a^n = r^n \cdot 0 = 0 \implies ra \in \text{Nil}(R)$. Ч.т.д.
%──────────────────────────────────────────────────────────────────────────────
\section{Главные идеалы. Является ли подкольцо идеалом?}

% Вопрос 17. Главные идеалы коммутативного кольца. Примеры...

\textbf{Определение.} Идеал $I$ коммутативного кольца $R$ с единицей называется главным, если он порождён одним элементом $a \in R$. Обозначается $I = (a) = \{ ra \mid r \in R \}$.

\textbf{Примеры:} В кольце $\mathbb{Z}$ любой идеал является главным и имеет вид $n\mathbb{Z} = (n)$. В кольце многочленов $\mathbb{R}[x]$ множество многочленов, делящихся на $x-1$, образует главный идеал $(x-1)$.

\textbf{Верно ли, что подкольцо является идеалом?} 
Нет, в общем случае это неверно. 
\textbf{Пример:} $\mathbb{Z}$ является подкольцом поля $\mathbb{Q}$ (замкнуто относительно сложения, вычитания, умножения). Однако $\mathbb{Z}$ \textbf{не является} идеалом в $\mathbb{Q}$. Возьмём $1 \in \mathbb{Z}$ и внешний элемент $\frac{1}{2} \in \mathbb{Q}$. Их произведение $\frac{1}{2} \cdot 1 = \frac{1}{2} \notin \mathbb{Z}$. Идеал обязан "поглощать" умножение на элементы объемлющего кольца.

\textbf{При каких условиях это верно?} 
Подкольцо $S$ кольца $R$ является идеалом, если выполняется условие поглощения: для любого элемента $s \in S$ и \textbf{любого} элемента $r \in R$ их произведения $rs$ и $sr$ принадлежат $S$ (то есть $R \cdot S \subseteq S$ и $S \cdot R \subseteq S$).

%──────────────────────────────────────────────────────────────────────────────
\section{Поле является простым кольцом, и наоборот}

\begin{definition}
	Кольцо $R$ называется \textbf{простым}, если оно ненулевое и не имеет нетривиальных двусторонних идеалов (то есть единственные идеалы --- $\{0\}$ и $R$).
\end{definition}

\begin{theorem}
	Поле является простым кольцом.
\end{theorem}
\begin{proof}
	Пусть $I \neq \{0\}$ --- идеал поля $F$, $a \in I$, $a \neq 0$. Тогда $a$ обратим и $e = a \cdot a^{-1} \in I$, откуда $I = F$. \qed
\end{proof}

\begin{theorem}
	Коммутативное кольцо с единицей является полем тогда и только тогда, когда оно простое.
\end{theorem}
\begin{proof}
	($\Rightarrow$) Доказано выше.
	
	($\Leftarrow$) Пусть $R$ --- коммутативное простое кольцо с единицей, $a \neq 0$. Тогда $Ra = \{ra \mid r \in R\}$ --- ненулевой идеал, значит $Ra = R$. Следовательно, существует $b$ с $ba = e$, то есть $a$ обратим. \qed
\end{proof}

%──────────────────────────────────────────────────────────────────────────────
\section{Кольцо матриц над полем является простым}

\begin{theorem}
	Кольцо $M_n(F)$ квадратных матриц $n \times n$ над полем $F$ является простым.
\end{theorem}
\begin{proof}
	Пусть $I \neq \{0\}$ --- двусторонний идеал $M_n(F)$, $A \in I$, $A \neq 0$. Тогда существуют индексы $i, j$ с $A_{ij} \neq 0$. Обозначим через $E_{kl}$ матрицу, у которой на месте $(k,l)$ стоит $1$, остальные элементы нулевые. Тогда:
	\[
	E_{ki} \cdot A \cdot E_{jl} = A_{ij} \cdot E_{kl}.
	\]
	Так как $I$ --- двусторонний идеал и $A \in I$, то $A_{ij} E_{kl} \in I$. Поскольку $A_{ij} \neq 0 \in F$, то $E_{kl} = A_{ij}^{-1} \cdot (A_{ij} E_{kl}) \in I$. Следовательно, все $E_{kl} \in I$ для всех $k, l$, откуда $I = M_n(F)$. \qed
\end{proof}

%──────────────────────────────────────────────────────────────────────────────
\section{Произведение идеалов является идеалом. Примеры}

% Вопрос 20. Произведение идеалов.

\textbf{Определение.} Произведением идеалов $I$ и $J$ кольца $R$ называется множество $IJ$, состоящее из всевозможных конечных сумм вида:
$$ IJ = \left\{ \sum_{k=1}^n i_k j_k \mid i_k \in I, j_k \in J, n \in \mathbb{N} \right\} $$
Просто множество произведений $\{ij\}$ не обязано быть замкнутым по сложению, поэтому берутся суммы.

\textbf{Доказательство, что $IJ$ — идеал.}
1. Замкнутость по разности: Разность двух конечных сумм такого вида — это снова конечная сумма такого вида, значит $IJ$ — подгруппа по сложению.
2. Поглощение умножения: Пусть $x = \sum i_k j_k \in IJ$ и $r \in R$. Тогда:
$rx = r \left( \sum i_k j_k \right) = \sum (r i_k) j_k$.
Так как $I$ — идеал, $(r i_k) = i'_k \in I$. Значит $rx = \sum i'_k j_k \in IJ$. Аналогично для правого умножения $xr \in IJ$. Ч.т.д.

\textbf{Примеры:}
1. \textit{В целых числах $\mathbb{Z}$:} Пусть $I = (4)$, $J = (6)$. Произведение порождённых ими главных идеалов равно идеалу, порожденному произведением образующих: $IJ = (4 \cdot 6) = (24)$.
2. \textit{В кольце многочленов $\mathbb{R}[x]$:} Пусть $I = (x^2)$ (многочлены без свободного члена и без $x$) и $J = (x-1)$ (многочлены с корнем 1). Их произведение $IJ = (x^2(x-1)) = (x^3 - x^2)$.

%──────────────────────────────────────────────────────────────────────────────
\section{Сумма идеалов. Дистрибутивность умножения относительно сложения}

% Вопрос 21. Сумма идеалов.

\textbf{Определение.} Суммой идеалов $I$ и $J$ называется множество $I + J = \{i + j \mid i \in I, j \in J\}$. Это наименьший идеал, содержащий и $I$, и $J$.

\textbf{Свойства:}
1. Коммутативность: $I + J = J + I$.
2. Ассоциативность: $(I + J) + K = I + (J + K)$.
3. Идемпотентность: $I + I = I$.
4. Если $I \subseteq J$, то $I + J = J$.

\textbf{Примеры:}
1. \textit{В $\mathbb{Z}$:} Сумма идеалов $(a) + (b) = (\text{НОД}(a, b))$. Например, $(4) + (6) = (2)$.
2. \textit{В $\mathbb{R}[x]$:} Идеалы $I = (x)$ и $J = (x-1)$. Так как $\text{НОД}(x, x-1) = 1$, их сумма равна $(1) = \mathbb{R}[x]$.

\textbf{Доказательство дистрибутивности: $I(J+K) = IJ + IK$}.
1) $\subseteq$: Любой элемент из $I(J+K)$ есть сумма вида $\sum i_m (j_m + k_m) = \sum i_m j_m + \sum i_m k_m$. Первая сумма принадлежит $IJ$, вторая $IK$, значит весь элемент в $IJ + IK$.
2) $\supseteq$: Любой элемент из $IJ + IK$ имеет вид $\sum i_m j_m + \sum i'_n k_n$. Каждое слагаемое $i_m j_m = i_m(j_m + 0) \in I(J+K)$ и $i'_n k_n = i'_n(0 + k_n) \in I(J+K)$. Значит, в силу замкнутости идеала $I(J+K)$ по сложению, и вся сумма принадлежит $I(J+K)$. Ч.т.д.

%──────────────────────────────────────────────────────────────────────────────
\section{Пересечение идеалов}

% Вопрос 22. Пересечение идеалов.

\textbf{Теорема.} Пересечение любого семейства идеалов $\bigcap_{a \in A} I_a$ является идеалом.
\textbf{Доказательство.} Пусть $x, y \in \bigcap I_a$. Это значит, что $\forall a: x, y \in I_a$. 
Так как каждый $I_a$ — идеал, разность $x - y \in I_a$ для любого $a$. Значит, $x - y \in \bigcap I_a$.
Пусть $r \in R$. Так как $x \in I_a$ и $I_a$ идеал, то $rx \in I_a$ и $xr \in I_a$ для всех $a$. Значит, $rx, xr \in \bigcap I_a$. Ч.т.д.

\textbf{Образует ли множество идеалов кольца кольцо?}
Нет. Операции сложения (сумма идеалов) и умножения идеалов определены и даже обладают нужной дистрибутивностью. Однако по сложению (которое играет роль суммы идеалов) \textbf{нет обратных элементов}. Для идеала $I \neq \{0\}$ не существует идеала $J$ такого, что $I + J = \{0\}$, так как $\{0\} \subset I \subset I+J$. Поэтому множество идеалов образует лишь полукольцо.

\textbf{Примеры:}
1. \textit{В $\mathbb{Z}$:} Пересечение $(a) \cap (b) = (\text{НОК}(a, b))$. Например, $(4) \cap (6) = (12)$.
2. \textit{В $\mathbb{R}[x]$:} Пересечение $(x) \cap (x-1)$ состоит из многочленов, имеющих корни 0 и 1, то есть порождается их произведением (т.к. они взаимно просты): $(x(x-1)) = (x^2 - x)$.

\textbf{Доказательство, что $IJ \subseteq I \cap J$:}
Пусть элемент $x \in IJ$. Тогда $x = \sum i_k j_k$.
Рассмотрим слагаемое $i_k j_k$. Так как $i_k \in I$ и $I$ идеал, произведение $i_k j_k \in I$. С другой стороны, $j_k \in J$ и $J$ идеал, значит $i_k j_k \in J$.
Таким образом, каждое слагаемое принадлежит и $I$, и $J$, а значит, принадлежит их пересечению $I \cap J$. Сумма таких слагаемых также принадлежит $I \cap J$. Следовательно, $IJ \subseteq I \cap J$.

%──────────────────────────────────────────────────────────────────────────────
\section{Взаимно простые идеалы}

\begin{definition}
	Идеалы $A,B$ коммутативного кольца с единицей называются \textbf{взаимно простыми}, если
	\[
	A+B=R.
	\]
\end{definition}

\begin{example}
	В кольце $\Z$ идеалы $(m)$ и $(n)$ взаимно просты тогда и только тогда, когда $(m,n)=1$.
\end{example}

\begin{theorem}
	Идеалы $A,B$ взаимно просты тогда и только тогда, когда существуют элементы
	\[
	a\in A,\qquad b\in B,
	\]
	такие, что
	\[
	a+b=e.
	\]
\end{theorem}

\begin{proof}
	Если $A+B=R$, то $e\in A+B$, следовательно,
	\[
	e=a+b
	\]
	для некоторых $a\in A$, $b\in B$.
	
	Обратно, если существует разложение $e=a+b$, то для любого $r\in R$
	\[
	r=re=r(a+b)=ra+rb.
	\]
	Так как $ra\in A$ и $rb\in B$, то $r\in A+B$. Следовательно,
	\[
	A+B=R.
	\]
\end{proof}

%──────────────────────────────────────────────────────────────────────────────
\section{Идеал и произведение взаимно простых идеалов}

% Вопрос 24. Взаимно простые идеалы с произведением. Идеал, порождённый конечным множеством.

\textbf{Взаимно простые идеалы:} Идеалы $I, J$ называются взаимно простыми, если $I + J = R$. Это равносильно тому, что $\exists i \in I, j \in J : i + j = 1$.

\textbf{Теорема.} Если идеал $I$ взаимно прост с идеалами $J_1$ и $J_2$, то он взаимно прост с их произведением $J_1 J_2$.
\textbf{Доказательство.} По условию $I + J_1 = R$ и $I + J_2 = R$. Следовательно, существуют элементы $i_1, i_2 \in I$, $j_1 \in J_1$ и $j_2 \in J_2$ такие, что:
$i_1 + j_1 = 1$ и $i_2 + j_2 = 1$.
Перемножим эти два равенства:
$$ 1 = (i_1 + j_1)(i_2 + j_2) = i_1 i_2 + i_1 j_2 + j_1 i_2 + j_1 j_2 $$
Первые три слагаемых принадлежат $I$ (так как $i_1, i_2 \in I$, а $I$ — идеал, выдерживающий умножение на элементы кольца). 
Последнее слагаемое $j_1 j_2 \in J_1 J_2$ по определению произведения идеалов.
Обозначим $i' = i_1 i_2 + i_1 j_2 + j_1 i_2 \in I$. Тогда $1 = i' + j_1 j_2 \in I + J_1 J_2$.
Раз единица принадлежит $I + J_1 J_2$, этот идеал совпадает со всем кольцом $R$, что и означает взаимную простоту $I$ и $J_1 J_2$. По индукции обобщается на любое конечное число идеалов.

\textbf{Идеал, порождённый конечным множеством элементов.} 
Если $X = \{x_1, \dots, x_n\}$ — конечное подмножество коммутативного кольца $R$ с единицей, то идеал, порожденный $X$, обозначается $(x_1, \dots, x_n)$ и состоит из всех возможных линейных комбинаций этих элементов с коэффициентами из $R$.
\textbf{Вид элементов:} Любой элемент $a \in (x_1, \dots, x_n)$ может быть представлен в виде:
$$ a = r_1 x_1 + r_2 x_2 + \dots + r_n x_n, \quad \text{где } r_i \in R $$
Если в кольце нет единицы, к этой сумме нужно добавить целочисленные кратные $\sum m_i x_i$, где $m_i \in \mathbb{Z}$.
%──────────────────────────────────────────────────────────────────────────────
\section{Простые и максимальные идеалы}

\begin{definition}
	Идеал $P\neq R$ коммутативного кольца называется \textbf{простым}, если из
	\[
	ab\in P
	\]
	следует
	\[
	a\in P \quad \text{или}\quad b\in P.
	\]
	Идеал $M\neq R$ называется \textbf{максимальным}, если между $M$ и $R$ нет промежуточных идеалов.
\end{definition}

\begin{theorem}
	Нулевой идеал коммутативного кольца с единицей прост тогда и только тогда, когда кольцо является областью целостности.
\end{theorem}

\begin{proof}
	Идеал $(0)$ прост тогда и только тогда, когда из
	\[
	ab=0
	\]
	следует
	\[
	a=0 \quad \text{или}\quad b=0.
	\]
	Это и есть определение области целостности.
\end{proof}

%──────────────────────────────────────────────────────────────────────────────
\section{Максимальный идеал прост}

\begin{theorem}
	Каждый максимальный идеал коммутативного кольца с единицей является простым.
\end{theorem}

\begin{proof}
	Пусть $M$ --- максимальный идеал и
	\[
	ab\in M,\qquad a\notin M.
	\]
	Рассмотрим идеал
	\[
	M+(a).
	\]
	Так как $M$ максимален и $a\notin M$, то
	\[
	M+(a)=R.
	\]
	Следовательно, существуют $m\in M$, $r\in R$ такие, что
	\[
	m+ra=e.
	\]
	Умножим на $b$:
	\[
	mb+rab=b.
	\]
	Так как $mb\in M$ и $ab\in M$, то $rab\in M$. Значит,
	\[
	b\in M.
	\]
	Следовательно, $M$ прост.
\end{proof}

%──────────────────────────────────────────────────────────────────────────────
\section{Необратимые элементы и максимальные идеалы}

\begin{theorem}
	Любой необратимый элемент кольца с единицей содержится в некотором максимальном идеале.
\end{theorem}

\begin{proof}
	Пусть $a$ необратим. Тогда главный идеал $(a)\neq R$. Рассмотрим множество всех собственных идеалов, содержащих $(a)$. По лемме Цорна существует максимальный элемент $M$. Он является максимальным идеалом кольца и содержит $a$.
\end{proof}

\begin{remark}
	В кольце $\Z$ простыми идеалами являются:
	\[
	(0),\qquad (p),
	\]
	где $p$ --- простое число.
\end{remark}

%──────────────────────────────────────────────────────────────────────────────
\section{Радикал кольца}

\begin{definition}
	Радикалом коммутативного кольца называется множество
	\[
	\mathrm{Rad}(R)=\{a\in R\mid a^n=0 \text{ для некоторого } n\ge1\}.
	\]
\end{definition}

\begin{theorem}
	Радикал кольца совпадает с пересечением всех простых идеалов:
	\[
	\mathrm{Rad}(R)=\bigcap_{P\text{ --- простой}} P.
	\]
\end{theorem}

\begin{proof}
	Если $a^n=0$, то $0\in P$ для любого простого идеала $P$. Из простоты $P$ следует $a\in P$. Значит,
	\[
	a\in \bigcap P.
	\]
	
	Обратно, если $a\notin \mathrm{Rad}(R)$, то множество
	\[
	S=\{1,a,a^2,\dots\}
	\]
	не содержит нуля. Существует простой идеал $P$, не пересекающийся с $S$. Тогда $a\notin P$, следовательно,
	\[
	a\notin \bigcap P.
	\]
\end{proof}

%──────────────────────────────────────────────────────────────────────────────
\section{Свойство простых идеалов}

\textbf{Утверждение 1.} Если простой идеал $P$ в коммутативном кольце $R$ содержит пересечение идеалов $A_1, A_2, \dots, A_n$, то один из них содержится в $P$.

\textbf{Доказательство.} Проведём доказательство от противного для двух идеалов $A$ и $B$. Пусть $A \cap B \subseteq P$, но при этом ни $A$, ни $B$ не содержатся в $P$. 
Это значит, что существуют элементы $a \in A \setminus P$ и $b \in B \setminus P$.
Рассмотрим их произведение $ab$. 
Так как $A$ и $B$ — идеалы, $ab \in A$ (т.к. $a \in A$) и $ab \in B$ (т.к. $b \in B$). 
Следовательно, $ab \in A \cap B$.
По условию $A \cap B \subseteq P$, значит, $ab \in P$.
Так как $P$ — простой идеал, из $ab \in P$ следует, что либо $a \in P$, либо $b \in P$. 
Но мы выбрали $a \notin P$ и $b \notin P$. Противоречие. Следовательно, либо $A \subseteq P$, либо $B \subseteq P$. 
По индукции это свойство обобщается на любое конечное число идеалов.

\textbf{Утверждение 2.} Если простой идеал $P$ совпадает с пересечением идеалов $A_1, \dots, A_n$, то он совпадает с одним из них.

\textbf{Доказательство.} Дано $P = \bigcap_{i=1}^n A_i$. 
По определению пересечения множеств, $P \subseteq A_i$ для всех $i = 1, \dots, n$.
С другой стороны, так как $P$ содержит пересечение идеалов, по Утверждению 1 существует такой индекс $k$, что $A_k \subseteq P$.
Объединяя два включения $P \subseteq A_k$ и $A_k \subseteq P$, получаем $P = A_k$. Ч.т.д.

%──────────────────────────────────────────────────────────────────────────────
\section{Факториальные кольца}

\begin{definition}
	Область целостности называется \textbf{факториальной}, если любой ненулевой необратимый элемент раскладывается в произведение неприводимых, причём это разложение единственно с точностью до порядка и ассоциированности множителей.
\end{definition}

\begin{definition}
	Ненулевой необратимый элемент $p$ называется \textbf{простым}, если из
	\[
	p\mid ab
	\]
	следует
	\[
	p\mid a \quad \text{или}\quad p\mid b.
	\]
\end{definition}

\begin{example}
	Факториальными являются кольца
	\[
	\Z,\qquad F[x],
	\]
	где $F$ --- поле.
\end{example}

\begin{remark}
	Не всякая область целостности является факториальной. Например,
	\[
	\Z[\sqrt{-5}]
	\]
	не факториально, так как
	\[
	6=2\cdot3=(1+\sqrt{-5})(1-\sqrt{-5})
	\]
	--- два различных разложения на неприводимые множители.
\end{remark}

\begin{theorem}
	В факториальном кольце существуют НОД и НОК любых ненулевых элементов.
\end{theorem}

\begin{proof}
	Разложим элементы на простые множители:
	\[
	a=u p_1^{\alpha_1}\cdots p_n^{\alpha_n},\qquad
	b=v p_1^{\beta_1}\cdots p_n^{\beta_n}.
	\]
	Тогда
	\[
	(a,b)=p_1^{\min(\alpha_1,\beta_1)}\cdots p_n^{\min(\alpha_n,\beta_n)},
	\]
	а
	\[
	[a,b]=p_1^{\max(\alpha_1,\beta_1)}\cdots p_n^{\max(\alpha_n,\beta_n)}.
	\]
\end{proof}

%──────────────────────────────────────────────────────────────────────────────
\section{Свойства взаимной простоты в факториальном кольце}

\begin{theorem}
	Если элемент $A$ взаимно прост с элементами $B,C,\dots,D$, то он взаимно прост с их произведением.
\end{theorem}

\begin{proof}
	Пусть некоторый простой делитель $p$ делит одновременно $A$ и произведение
	\[
	BCD\cdots D.
	\]
	Тогда $p$ делит один из множителей, например $B$. Получаем, что $p$ делит и $A$, и $B$, противоречие.
\end{proof}

\begin{theorem}
	Если элементы $B,C,\dots,D$ попарно взаимно просты и элемент $A$ делится на каждый из них, то
	\[
	A \text{ делится на } BCD\cdots D.
	\]
\end{theorem}

\begin{proof}
	Так как множители попарно взаимно просты, их разложения на простые множители не пересекаются. Поскольку каждый из них делит $A$, все соответствующие простые множители входят в разложение $A$. Следовательно,
	\[
	BCD\cdots D\mid A.
	\]
\end{proof}

%──────────────────────────────────────────────────────────────────────────────
\section{Свойства простых элементов в факториальном кольце}

\begin{theorem}
	Пусть $p$ --- простой элемент факториального кольца $R$. Тогда для любого $a\in R$:
	\[
	p\mid a \quad \text{или}\quad (p,a)=1.
	\]
\end{theorem}

\begin{proof}
	Если $p\nmid a$, то в разложении $a$ на простые множители элемент $p$ не встречается. Следовательно, общих простых делителей у $p$ и $a$ нет, то есть
	\[
	(p,a)=1.
	\]
\end{proof}

\begin{theorem}
	Если
	\[
	p\mid abc\cdots d,
	\]
	то хотя бы один из множителей делится на $p$.
\end{theorem}

\begin{proof}
	Так как $p$ прост, из
	\[
	p\mid ab
	\]
	следует
	\[
	p\mid a \quad \text{или}\quad p\mid b.
	\]
	Применяя это свойство последовательно к произведению нескольких множителей, получаем утверждение теоремы.
\end{proof}

%──────────────────────────────────────────────────────────────────────────────
\section{НОД в кольце главных идеалов}

\begin{theorem}
	В кольце главных идеалов любые два элемента имеют НОД.
\end{theorem}

\begin{proof}
	Пусть
	\[
	a,b\in R.
	\]
	Рассмотрим идеал
	\[
	(a)+(b).
	\]
	Так как кольцо является кольцом главных идеалов,
	\[
	(a)+(b)=(d)
	\]
	для некоторого элемента $d$.
	
	Так как
	\[
	a,b\in(d),
	\]
	то
	\[
	d\mid a,
	\qquad
	d\mid b.
	\]
	Следовательно, $d$ --- общий делитель.
	
	Если $c$ --- любой общий делитель элементов $a,b$, то
	\[
	a,b\in(c),
	\]
	значит,
	\[
	(a)+(b)\subseteq(c).
	\]
	Отсюда
	\[
	d\in(c),
	\]
	то есть
	\[
	c\mid d.
	\]
	Следовательно, $d$ является НОД.
\end{proof}

\begin{theorem}
	В кольце главных идеалов НОД представим в виде:
	\[
	d=ax+by.
	\]
\end{theorem}

\begin{proof}
	Так как
	\[
	(a)+(b)=(d),
	\]
	то
	\[
	d\in(a)+(b).
	\]
	Следовательно,
	\[
	d=ax+by
	\]
	для некоторых
	\[
	x,y\in R.
	\]
\end{proof}

%──────────────────────────────────────────────────────────────────────────────
\section{Свойства неприводимых элементов в КГИ}

\begin{theorem}
	Пусть $p$ --- неприводимый элемент кольца главных идеалов. Тогда для любого
	$a\in R$:
	\[
	p\mid a
	\quad \text{или}\quad
	(p,a)=1.
	\]
\end{theorem}

\begin{proof}
	Рассмотрим НОД:
	\[
	d=(p,a).
	\]
	Так как
	\[
	d\mid p,
	\]
	а $p$ неприводим, то
	\[
	d\sim p
	\quad \text{или}\quad
	d\sim 1.
	\]
	В первом случае:
	\[
	p\mid a.
	\]
	Во втором:
	\[
	(p,a)=1.
	\]
\end{proof}

\begin{theorem}
	Если
	\[
	p\mid a_1a_2\cdots a_n,
	\]
	то хотя бы один из множителей делится на $p$.
\end{theorem}

\begin{proof}
	Достаточно доказать для двух множителей.
	
	Пусть
	\[
	p\mid ab,
	\qquad
	p\nmid a.
	\]
	Тогда
	\[
	(p,a)=1.
	\]
	Следовательно, существуют
	\[
	x,y\in R,
	\]
	такие, что
	\[
	px+ay=1.
	\]
	Умножая на $b$:
	\[
	pb+aby=b.
	\]
	Так как
	\[
	p\mid pb,
	\qquad
	p\mid ab,
	\]
	то
	\[
	p\mid b.
	\]
\end{proof}

%──────────────────────────────────────────────────────────────────────────────
\section{Неразложимые элементы}

\begin{theorem}
	Пусть $a$ --- ненулевой необратимый элемент области целостности, который не
	разлагается на неприводимые множители. Тогда $a$ имеет собственный делитель,
	который также не разлагается на неприводимые множители.
\end{theorem}

\begin{proof}
	Так как $a$ не является неприводимым, существуют необратимые элементы
	\[
	b,c
	\]
	такие, что
	\[
	a=bc.
	\]
	
	Если бы оба элемента раскладывались на неприводимые множители, то и $a$
	разлагался бы на неприводимые множители. Противоречие.
	
	Следовательно, хотя бы один из элементов $b,c$ не раскладывается на
	неприводимые множители. Он является собственным делителем элемента $a$.
\end{proof}

%──────────────────────────────────────────────────────────────────────────────
\section{Бесконечные цепочки делителей и идеалов}

\begin{theorem}
	Если в области целостности существует ненулевой необратимый элемент,
	не разлагающийся на неприводимые множители, то существует бесконечная
	последовательность:
	\[
	a_1,a_2,a_3,\dots
	\]
	такая, что каждый следующий элемент является собственным делителем
	предыдущего.
\end{theorem}

\begin{proof}
	Пусть
	\[
	a_1=a.
	\]
	По предыдущей теореме существует собственный делитель
	\[
	a_2\mid a_1,
	\]
	не разлагающийся на неприводимые множители.
	
	Продолжая процесс, получаем последовательность:
	\[
	a_1,a_2,a_3,\dots
	\]
	где
	\[
	a_{n+1}\mid a_n,
	\]
	и деление является собственным.
\end{proof}

\begin{theorem}
	Тогда существует бесконечная возрастающая цепочка главных идеалов:
	\[
	(a_1)\subset (a_2)\subset (a_3)\subset\dots
	\]
\end{theorem}

\begin{proof}
	Так как
	\[
	a_{n+1}\mid a_n,
	\]
	то
	\[
	(a_n)\subseteq(a_{n+1}).
	\]
	Так как делитель собственный, включение строгое.
\end{proof}

%──────────────────────────────────────────────────────────────────────────────
\section{Стабилизация идеалов в КГИ}

\begin{theorem}
	В кольце главных идеалов любая возрастающая цепочка идеалов стабилизируется.
\end{theorem}

\begin{proof}
	Пусть
	\[
	I_1\subseteq I_2\subseteq I_3\subseteq\dots
	\]
	--- возрастающая цепочка идеалов.
	
	Рассмотрим объединение:
	\[
	I=\bigcup_{n=1}^{\infty} I_n.
	\]
	Легко проверить, что $I$ является идеалом.
	
	Так как кольцо является КГИ,
	\[
	I=(a).
	\]
	Элемент
	\[
	a\in I_n
	\]
	для некоторого $n$. Тогда
	\[
	I=(a)\subseteq I_n.
	\]
	Но также
	\[
	I_n\subseteq I.
	\]
	Следовательно,
	\[
	I_n=I.
	\]
	Значит, для всех
	\[
	m\ge n
	\]
	имеем
	\[
	I_m=I_n.
	\]
\end{proof}

\begin{theorem}
	Каждый ненулевой необратимый элемент КГИ раскладывается на неприводимые
	множители.
\end{theorem}

\begin{proof}
	Если бы существовал элемент без такого разложения, то по предыдущей теореме
	существовала бы бесконечная строго возрастающая цепочка главных идеалов,
	что невозможно.
\end{proof}

%──────────────────────────────────────────────────────────────────────────────
\section{Факториальность КГИ}

\begin{theorem}
	Любое кольцо главных идеалов является факториальным.
\end{theorem}

\begin{proof}
	Существование разложения уже доказано.
	
	Докажем единственность. Пусть
	\[
	p
	\]
	--- неприводимый элемент. Тогда из
	\[
	p\mid ab
	\]
	следует
	\[
	p\mid a
	\quad \text{или}\quad
	p\mid b.
	\]
	Следовательно, неприводимые элементы являются простыми.
	
	Теперь стандартным аргументом доказывается единственность разложения на
	неприводимые множители с точностью до порядка и ассоциированности.
\end{proof}

%──────────────────────────────────────────────────────────────────────────────
\section{Евклидово кольцо и КГИ}

\begin{theorem}
	Любое евклидово кольцо является кольцом главных идеалов.
\end{theorem}

\begin{proof}
	Пусть
	\[
	I
	\]
	--- ненулевой идеал евклидова кольца.
	
	Выберем ненулевой элемент
	\[
	d\in I
	\]
	с минимальным значением евклидовой нормы.
	
	Докажем:
	\[
	I=(d).
	\]
	
	Для любого
	\[
	a\in I
	\]
	разделим с остатком:
	\[
	a=qd+r,
	\]
	где
	\[
	r=0
	\]
	или
	\[
	\nu(r)<\nu(d).
	\]
	
	Но
	\[
	r=a-qd\in I.
	\]
	По минимальности нормы:
	\[
	r=0.
	\]
	Следовательно,
	\[
	a\in(d).
	\]
	
	Значит,
	\[
	I=(d).
	\]
\end{proof}

%──────────────────────────────────────────────────────────────────────────────
\section{Простые идеалы в КГИ}

\begin{theorem}
	В кольце главных идеалов каждый простой идеал максимален.
\end{theorem}

\begin{proof}
	Пусть
	\[
	P
	\]
	--- простой идеал.
	
	Так как кольцо является КГИ,
	\[
	P=(p).
	\]
	
	Рассмотрим идеал
	\[
	I,
	\qquad
	P\subseteq I\subseteq R.
	\]
	Тогда
	\[
	I=(a)
	\]
	для некоторого $a$.
	
	Из
	\[
	(p)\subseteq(a)
	\]
	следует
	\[
	a\mid p.
	\]
	
	Так как $p$ прост и, следовательно, неприводим:
	\[
	a\sim p
	\quad \text{или}\quad
	a\sim 1.
	\]
	
	В первом случае:
	\[
	I=P.
	\]
	Во втором:
	\[
	I=R.
	\]
	
	Следовательно, промежуточных идеалов нет, и $P$ максимален.
\end{proof}

%──────────────────────────────────────────────────────────────────────────────
\section{Тела и простые кольца}

\begin{definition}
	Тело --- это кольцо с единицей, в котором каждый ненулевой элемент обратим.
\end{definition}

\begin{example}
	Тело кватернионов:
	\[
	\mathbb{H}
	\]
	является некоммутативным телом.
\end{example}

\begin{theorem}
	Любое тело является простым кольцом.
\end{theorem}

\begin{proof}
	Пусть
	\[
	I
	\]
	--- ненулевой идеал тела $D$.
	
	Возьмём
	\[
	0\neq a\in I.
	\]
	Так как $a$ обратим,
	\[
	a^{-1}a=1\in I.
	\]
	Тогда для любого
	\[
	x\in D
	\]
	имеем
	\[
	x=x\cdot1\in I.
	\]
	Следовательно,
	\[
	I=D.
	\]
	
	Таким образом, у тела нет нетривиальных идеалов.
\end{proof}

%──────────────────────────────────────────────────────────────────────────────
\section{Кольцо без нетривиальных левых идеалов}

\begin{theorem}
	Ненулевое кольцо с единицей, не имеющее нетривиальных левых идеалов,
	является телом.
\end{theorem}

\begin{proof}
	Пусть
	\[
	0\neq a\in R.
	\]
	Рассмотрим левый идеал:
	\[
	Ra=\{xa\mid x\in R\}.
	\]
	
	Так как
	\[
	a\neq0,
	\]
	идеал ненулевой. По условию:
	\[
	Ra=R.
	\]
	
	Следовательно,
	\[
	1\in Ra,
	\]
	то есть существует
	\[
	x\in R
	\]
	такой, что
	\[
	xa=1.
	\]
	
	Аналогично существует
	\[
	y\in R
	\]
	такой, что
	\[
	ay=1.
	\]
	
	Тогда
	\[
	x=x(ay)=(xa)y=y.
	\]
	Следовательно,
	\[
	xa=ax=1.
	\]
	
	Значит, любой ненулевой элемент обратим, и $R$ является телом.
\end{proof}

%──────────────────────────────────────────────────────────────────────────────
\section{Центр кольца}

\begin{definition}
	Центром кольца $R$ называется множество:
	\[
	Z(R)=\{a\in R\mid ax=xa \text{ для всех } x\in R\}.
	\]
\end{definition}

\begin{theorem}
	Центр кольца является коммутативным подкольцом.
\end{theorem}

\begin{proof}
	Если
	\[
	a,b\in Z(R),
	\]
	то для любого
	\[
	x\in R
	\]
	имеем:
	\[
	(a+b)x=ax+bx=xa+xb=x(a+b),
	\]
	и
	\[
	(ab)x=a(bx)=a(xb)=(ax)b=(xa)b=x(ab).
	\]
	
	Следовательно,
	\[
	a+b,ab\in Z(R).
	\]
	Также
	\[
	ab=ba,
	\]
	поэтому центр коммутативен.
\end{proof}

%──────────────────────────────────────────────────────────────────────────────
\section{Центр кольца и тела}

\begin{definition}
	Центром кольца $R$ называется множество
	\[
	Z(R)=\{a\in R\mid ax=xa \text{ для всех } x\in R\}.
	\]
\end{definition}

\begin{theorem}
	Центр кольца является коммутативным подкольцом кольца $R$.
\end{theorem}

\begin{proof}
	Если $a,b\in Z(R)$, то для любого $x\in R$
	\[
	(a+b)x=ax+bx=xa+xb=x(a+b),
	\]
	следовательно,
	\[
	a+b\in Z(R).
	\]
	
	Аналогично
	\[
	(ab)x=a(bx)=a(xb)=(ax)b=(xa)b=x(ab),
	\]
	поэтому
	\[
	ab\in Z(R).
	\]
	
	Также
	\[
	(-a)x=-(ax)=-(xa)=x(-a),
	\]
	значит,
	\[
	-a\in Z(R).
	\]
\end{proof}

\begin{example}
	\[
	Z(M_n(F))=\{\lambda E_n\mid \lambda\in F\}.
	\]
\end{example}

%──────────────────────────────────────────────────────────────────────────────
\section{Линейная структура тела над подтелом}

% Вопрос 45. Сравнимые элементы по модулю идеала...

\textbf{Определение:} Элементы $a, b \in R$ сравнимы по модулю идеала $I$ ($a \equiv b \pmod I$), если $a - b \in I$.

\textbf{Свойства эквивалентности:}
1. Рефлексивность: $a \equiv a \pmod I$, так как $a - a = 0 \in I$.
2. Симметричность: Если $a \equiv b \pmod I$, то $a - b \in I \implies -(a - b) = b - a \in I \implies b \equiv a \pmod I$.
3. Транзитивность: Если $a \equiv b$ и $b \equiv c$, то $(a - b) \in I$ и $(b - c) \in I \implies (a - b) + (b - c) = a - c \in I \implies a \equiv c \pmod I$.

\textbf{Случай нулевого идеала:} Если $I = \{0\}$, то $a \equiv b \pmod{\{0\}} \iff a - b = 0 \iff a = b$. Отношение вырождается в обычное равенство.

\textbf{Степени и кратные:} Пусть $a \equiv b \pmod I$ ($a - b \in I$).
1. Кратные: для любого $m \in \mathbb{Z}$, $ma - mb = m(a-b)$. Так как $a-b \in I$, а идеал является подгруппой, $m(a-b) \in I \implies ma \equiv mb \pmod I$.
2. Степени (в коммутативном кольце): $a^n - b^n = (a-b)(a^{n-1} + a^{n-2}b + \dots + b^{n-1})$. Так как $a-b \in I$ и $I$ выдерживает умножение на элементы кольца, произведение принадлежит $I \implies a^n \equiv b^n \pmod I$.

\textbf{Китайская теорема об остатках (КТО).} 
\textit{Для колец:} Если $I_1, \dots, I_k$ — попарно взаимно простые идеалы кольца $R$ с единицей ($I_i + I_j = R$ при $i \neq j$), то для любых $a_1, \dots, a_k \in R$ система сравнений $x \equiv a_i \pmod{I_i}$ имеет решение, причём единственное по модулю пересечения идеалов $\bigcap I_i$. Более того, $R / \bigcap I_i \cong \bigoplus (R / I_i)$.
\textit{Для целых чисел:} Если $m_1, \dots, m_k$ — попарно взаимно простые натуральные числа, то система $x \equiv a_i \pmod{m_i}$ имеет единственное решение по модулю $M = m_1 \dots m_k$.

%──────────────────────────────────────────────────────────────────────────────
\section{Сравнения по модулю идеала}

\begin{definition}
	Пусть $I$ --- идеал кольца $R$. Элементы $a,b\in R$ называются сравнимыми по модулю $I$, если
	\[
	a-b\in I.
	\]
	Обозначение:
	\[
	a\equiv b\pmod I.
	\]
\end{definition}

\begin{theorem}
	Отношение сравнимости по модулю идеала является отношением эквивалентности.
\end{theorem}

\begin{proof}
	\textbf{Рефлексивность:}
	\[
	a-a=0\in I.
	\]
	
	\textbf{Симметричность:}
	если
	\[
	a-b\in I,
	\]
	то
	\[
	b-a=-(a-b)\in I.
	\]
	
	\textbf{Транзитивность:}
	если
	\[
	a-b\in I,\qquad b-c\in I,
	\]
	то
	\[
	a-c=(a-b)+(b-c)\in I.
	\]
\end{proof}

\begin{remark}
	При $I=(0)$ сравнимость совпадает с равенством.
\end{remark}

\begin{definition}
	Классом вычетов элемента $a$ называется множество
	\[
	[a]=a+I=\{a+x\mid x\in I\}.
	\]
\end{definition}

\begin{theorem}
	Если
	\[
	a\equiv b\pmod I,
	\]
	то
	\[
	a^n\equiv b^n\pmod I,
	\]
	и для любого $r\in R$
	\[
	ra\equiv rb\pmod I.
	\]
\end{theorem}

\begin{proof}
	Так как
	\[
	a-b\in I,
	\]
	то
	\[
	a^n-b^n=(a-b)(a^{n-1}+a^{n-2}b+\dots+b^{n-1})\in I.
	\]
	Также
	\[
	ra-rb=r(a-b)\in I.
	\]
\end{proof}

\begin{theorem}[Китайская теорема об остатках]
	Пусть идеалы
	\[
	I_1,\dots,I_n
	\]
	попарно взаимно просты. Тогда отображение
	\[
	R\to R/I_1\times\dots\times R/I_n
	\]
	является сюръективным, а его ядро равно
	\[
	I_1\cap\dots\cap I_n.
	\]
\end{theorem}

%──────────────────────────────────────────────────────────────────────────────
\section{Простые и максимальные идеалы факторкольца}

\begin{theorem}
	Пусть $R$ --- коммутативное кольцо с единицей.
	
	Тогда идеал $P$ прост тогда и только тогда, когда
	\[
	R/P
	\]
	является областью целостности.
\end{theorem}

\begin{proof}
	В факторкольце
	\[
	(a+P)(b+P)=P
	\]
	тогда и только тогда, когда
	\[
	ab\in P.
	\]
	Так как $P$ прост,
	\[
	a\in P \quad \text{или}\quad b\in P.
	\]
	Следовательно,
	\[
	a+P=P \quad \text{или}\quad b+P=P.
	\]
	Значит, делителей нуля нет.
\end{proof}

\begin{theorem}
	Идеал $M$ максимален тогда и только тогда, когда
	\[
	R/M
	\]
	является полем.
\end{theorem}

\begin{proof}
	Пусть $M$ максимален и
	\[
	a+M\neq M.
	\]
	Тогда $a\notin M$. Из максимальности:
	\[
	M+(a)=R.
	\]
	Следовательно,
	\[
	m+ra=e
	\]
	для некоторых $m\in M$, $r\in R$. Тогда
	\[
	(r+M)(a+M)=e+M,
	\]
	то есть каждый ненулевой элемент обратим.
	
	Обратно, если $R/M$ --- поле, то факторкольцо не имеет нетривиальных идеалов. Следовательно, между $M$ и $R$ нет промежуточных идеалов, значит $M$ максимален.
\end{proof}

\begin{remark}
	Для кольца многочленов $F[x]$ над полем $F$ идеал $(f)$ максимален тогда и только тогда, когда многочлен $f$ неприводим.
\end{remark}

%──────────────────────────────────────────────────────────────────────────────
\section{Гомоморфизмы колец}

\begin{definition}
	Отображение
	\[
	f:R\to S
	\]
	называется гомоморфизмом колец, если
	\[
	f(a+b)=f(a)+f(b),
	\]
	\[
	f(ab)=f(a)f(b).
	\]
\end{definition}

\begin{theorem}
	Гомоморфный образ кольца является кольцом.
\end{theorem}

\begin{proof}
	Множество
	\[
	f(R)=\{f(a)\mid a\in R\}
	\]
	замкнуто относительно сложения и умножения:
	\[
	f(a)+f(b)=f(a+b),
	\]
	\[
	f(a)f(b)=f(ab).
	\]
	Также
	\[
	-f(a)=f(-a).
	\]
	Следовательно, $f(R)$ является кольцом.
\end{proof}

\begin{remark}
	Гомоморфный образ области целостности не обязан быть областью целостности.
	
	Например,
	\[
	\Z\to \Z/6\Z.
	\]
\end{remark}

\begin{remark}
	Гомоморфный образ кольца с делителями нуля может быть областью целостности.
	
	Например,
	\[
	\Z/6\Z \to \Z/3\Z.
	\]
\end{remark}

%──────────────────────────────────────────────────────────────────────────────
\section{Ядро гомоморфизма и теорема об изоморфизме}

\begin{definition}
	Ядром гомоморфизма
	\[
	f:R\to S
	\]
	называется множество
	\[
	\ker f=\{a\in R\mid f(a)=0\}.
	\]
\end{definition}

\begin{theorem}
	Ядро гомоморфизма является идеалом.
\end{theorem}

\begin{proof}
	Если
	\[
	a,b\in\ker f,
	\]
	то
	\[
	f(a-b)=f(a)-f(b)=0.
	\]
	Следовательно,
	\[
	a-b\in\ker f.
	\]
	
	Для любого $r\in R$
	\[
	f(ra)=f(r)f(a)=0,
	\]
	поэтому
	\[
	ra\in\ker f.
	\]
\end{proof}

\begin{theorem}[Основная теорема о гомоморфизме]
	Для любого гомоморфизма колец
	\[
	f:R\to S
	\]
	имеем изоморфизм
	\[
	R/\ker f \cong f(R).
	\]
\end{theorem}

\begin{proof}
	Определим отображение
	\[
	\varphi:R/\ker f\to f(R),
	\qquad
	\varphi(a+\ker f)=f(a).
	\]
	
	Оно корректно, так как из
	\[
	a-b\in\ker f
	\]
	следует
	\[
	f(a)=f(b).
	\]
	
	Отображение является сюръективным по определению образа. Ядро тривиально, следовательно, $\varphi$ инъективно.
\end{proof}

\begin{corollary}
	Если
	\[
	\ker f=(0),
	\]
	то гомоморфизм инъективен.
\end{corollary}

%──────────────────────────────────────────────────────────────────────────────
\section{Прообразы простых и максимальных идеалов}

% Вопрос 50. Характеристика области целостности. Бином Ньютона...

\textbf{Определение.} Характеристикой кольца $R$ с единицей называется наименьшее натуральное число $n > 0$ такое, что $n \cdot 1 = \underbrace{1 + \dots + 1}_{n} = 0$. Если такого $n$ не существует, говорят, что характеристика равна нулю ($\text{char } R = 0$). В области целостности характеристика может быть только нулём или простым числом $p$.

\textbf{Бином Ньютона в области целостности характеристики $p$:}
В области целостности простой характеристики $p$ формула бинома $(a+b)^p$ кардинально упрощается:
$$ (a+b)^p = a^p + b^p $$
\textbf{Доказательство:} Стандартная формула: $(a+b)^p = \sum_{k=0}^p \binom{p}{k} a^k b^{p-k}$. 
Биномиальный коэффициент $\binom{p}{k} = \frac{p!}{k!(p-k)!}$. Для $0 < k < p$ числитель делится на $p$, а знаменатель — нет (так как $p$ — простое число, а все множители знаменателя строго меньше $p$). Следовательно, $\binom{p}{k}$ делится на $p$.
В кольце характеристики $p$ любой элемент, умноженный на $p$, равен нулю: $p \cdot x = 0$. Значит, все промежуточные слагаемые зануляются, и остаётся только $a^p + b^p$.

\textbf{Простое поле.} Поле называется простым, если оно не содержит собственных подполей.

\textbf{Доказательство простоты полей $\mathbb{Z}_p$ и $\mathbb{Q}$:}
1. \textbf{Поле вычетов $\mathbb{Z}_p$:} Пусть $L \subseteq \mathbb{Z}_p$ — подполе. Оно должно содержать единицу $1$. В силу замкнутости относительно сложения, оно обязано содержать $1+1=2, 1+1+1=3, \dots, (p-1) \cdot 1 = p-1$. То есть $L$ содержит все элементы $\mathbb{Z}_p$, значит $L = \mathbb{Z}_p$.
2. \textbf{Поле рациональных чисел $\mathbb{Q}$:} Пусть $L \subseteq \mathbb{Q}$ — подполе. Аналогично, $1 \in L \implies \mathbb{Z} \subseteq L$. Так как $L$ замкнуто относительно деления (взятия обратного), то для любых $p \in \mathbb{Z}, q \in \mathbb{Z} \setminus \{0\}$, дробь $p q^{-1} = \frac{p}{q} \in L$. Следовательно, $L$ содержит все рациональные числа, $L = \mathbb{Q}$.

%──────────────────────────────────────────────────────────────────────────────
\section{Характеристика и простые поля}

\begin{definition}
	Характеристикой области целостности называется минимальное натуральное число $n$, такое что
	\[
	n\cdot 1=0.
	\]
	Если такого числа нет, характеристика равна нулю.
\end{definition}

\begin{theorem}
	Характеристика области целостности равна либо нулю, либо простому числу.
\end{theorem}

\begin{proof}
	Если
	\[
	\operatorname{char} R=ab,
	\qquad 1<a,b<n,
	\]
	то
	\[
	(a\cdot1)(b\cdot1)=0.
	\]
	Так как область целостности не имеет делителей нуля,
	\[
	a\cdot1=0
	\quad \text{или}\quad
	b\cdot1=0,
	\]
	противоречие минимальности.
\end{proof}

\begin{theorem}
	В области целостности характеристики $p$ выполняется формула:
	\[
	(a+b)^p=a^p+b^p.
	\]
\end{theorem}

\begin{proof}
	По биному Ньютона:
	\[
	(a+b)^p=\sum_{k=0}^p \binom{p}{k}a^k b^{p-k}.
	\]
	Для
	\[
	1\le k\le p-1
	\]
	число
	\[
	\binom{p}{k}
	\]
	делится на $p$, поэтому соответствующие слагаемые равны нулю.
\end{proof}

\begin{definition}
	Поле называется простым, если не содержит собственных подполей.
\end{definition}

\begin{theorem}
	Поля
	\[
	\Q
	\quad \text{и}\quad
	\mathbb{F}_p
	\]
	являются простыми.
\end{theorem}

\begin{proof}
	Любое подполе должно содержать единицу, а значит и все элементы вида
	\[
	n\cdot1.
	\]
	Следовательно, в характеристике $0$ подполе содержит $\Q$, а в характеристике $p$ содержит всё поле $\mathbb{F}_p$.
\end{proof}

%──────────────────────────────────────────────────────────────────────────────
\section{Полиномиальные функции и подстановка}

\begin{definition}
	Пусть
	\[
	f(x)\in R[x].
	\]
	Полиномиальной функцией называется отображение
	\[
	a\mapsto f(a).
	\]
\end{definition}

\begin{remark}
	В конечных кольцах ненулевой многочлен может задавать нулевую функцию.
	
	Например, над полем
	\[
	\mathbb{F}_p
	\]
	многочлен
	\[
	x^p-x
	\]
	обращается в нуль на всех элементах поля.
\end{remark}

\begin{theorem}
	Поле комплексных чисел изоморфно факторкольцу:
	\[
	\mathbb{C}\cong \mathbb{R}[x]/(x^2+1).
	\]
\end{theorem}

\begin{proof}
	Рассмотрим гомоморфизм
	\[
	\varphi:\mathbb{R}[x]\to\mathbb{C},
	\qquad
	\varphi(f)=f(i).
	\]
	
	Он сюръективен, поскольку
	\[
	a+bi=\varphi(a+bx).
	\]
	
	Ядро состоит из многочленов, делящихся на
	\[
	x^2+1,
	\]
	так как
	\[
	i^2+1=0.
	\]
	
	Следовательно,
	\[
	\ker\varphi=(x^2+1).
	\]
	
	По основной теореме о гомоморфизме:
	\[
	\mathbb{R}[x]/(x^2+1)\cong\mathbb{C}.
	\]
\end{proof}

%──────────────────────────────────────────────────────────────────────────────
\section{Алгебраические элементы}

\begin{definition}
	Элемент $\alpha$ расширения поля $L$ над полем $K$ называется алгебраическим над $K$, если существует ненулевой многочлен
	\[
	f(x)\in K[x]
	\]
	такой, что
	\[
	f(\alpha)=0.
	\]
\end{definition}

\begin{theorem}
	Для элемента $\alpha$ эквивалентны следующие условия:
	\begin{enumerate}
		\item $\alpha$ алгебраичен над $K$;
		\item множество
		\[
		1,\alpha,\alpha^2,\dots
		\]
		линейно зависимо над $K$;
		\item поле
		\[
		K(\alpha)
		\]
		имеет конечную степень над $K$.
	\end{enumerate}
\end{theorem}

\begin{proof}
	\[
	1\Rightarrow2
	\]
	Если
	\[
	a_0+a_1\alpha+\dots+a_n\alpha^n=0,
	\]
	то степени линейно зависимы.
	
	\[
	2\Rightarrow3
	\]
	Тогда все высокие степени выражаются через меньшие, поэтому
	\[
	K(\alpha)
	\]
	порождается конечным числом элементов.
	
	\[
	3\Rightarrow1
	\]
	Конечномерное пространство не может содержать бесконечно много линейно независимых степеней элемента.
\end{proof}

%──────────────────────────────────────────────────────────────────────────────
\section{Трансцендентные элементы}

\begin{definition}
	Элемент $\alpha$ расширения поля $L$ над полем $K$ называется
	\textbf{трансцендентным} над $K$, если не существует ненулевого многочлена
	\[
	f(x)\in K[x]
	\]
	такого, что
	\[
	f(\alpha)=0.
	\]
\end{definition}

\begin{theorem}
	Для элемента $\alpha$ эквивалентны следующие свойства:
	\begin{enumerate}
		\item $\alpha$ трансцендентен над $K$;
		\item система
		\[
		1,\alpha,\alpha^2,\dots
		\]
		линейно независима над $K$;
		\item поле
		\[
		K(\alpha)
		\]
		изоморфно полю рациональных дробей
		\[
		K(x).
		\]
	\end{enumerate}
\end{theorem}

\begin{proof}
	\[
	1\Rightarrow2
	\]
	Если степени линейно зависимы, то существует ненулевой многочлен,
	обращающийся в нуль на $\alpha$.
	
	\[
	2\Rightarrow3
	\]
	Отображение
	\[
	K[x]\to K[\alpha],
	\qquad
	f(x)\mapsto f(\alpha)
	\]
	инъективно, поэтому
	\[
	K[\alpha]\cong K[x].
	\]
	Тогда поле частных:
	\[
	K(\alpha)\cong K(x).
	\]
	
	\[
	3\Rightarrow1
	\]
	В поле $K(x)$ переменная $x$ не удовлетворяет никакому ненулевому
	многочлену над $K$.
\end{proof}

%──────────────────────────────────────────────────────────────────────────────
\section{Минимальный многочлен}

\begin{definition}
	Минимальным многочленом алгебраического элемента $\alpha$ над полем $K$
	называется неприводимый унитарный многочлен
	\[
	m_\alpha(x)\in K[x]
	\]
	минимальной степени, такой что
	\[
	m_\alpha(\alpha)=0.
	\]
\end{definition}

\begin{theorem}
	Если
	\[
	f(\alpha)=0,
	\]
	то минимальный многочлен делит $f$:
	\[
	m_\alpha(x)\mid f(x).
	\]
\end{theorem}

\begin{proof}
	Разделим $f$ с остатком:
	\[
	f=q m_\alpha+r,
	\qquad
	\deg r<\deg m_\alpha.
	\]
	Подставляя $\alpha$:
	\[
	0=f(\alpha)=q(\alpha)m_\alpha(\alpha)+r(\alpha)=r(\alpha).
	\]
	По минимальности степени:
	\[
	r=0.
	\]
	Следовательно,
	\[
	m_\alpha\mid f.
	\]
\end{proof}

\begin{definition}
	Степенью алгебраического элемента называется число
	\[
	[K(\alpha):K].
	\]
\end{definition}

\begin{theorem}
	Степень алгебраического элемента равна степени его минимального многочлена.
\end{theorem}

\begin{proof}
	Базис пространства
	\[
	K(\alpha)
	\]
	над $K$ имеет вид
	\[
	1,\alpha,\dots,\alpha^{n-1},
	\]
	где
	\[
	n=\deg m_\alpha.
	\]
	Следовательно,
	\[
	[K(\alpha):K]=n.
	\]
\end{proof}

\begin{remark}
	Элементы первой степени имеют вид:
	\[
	\alpha\in K.
	\]
\end{remark}

%──────────────────────────────────────────────────────────────────────────────
\section{Неприводимые многочлены}

\begin{theorem}
	Неприводимые многочлены над $\C$ имеют степень $1$.
\end{theorem}

\begin{proof}
	По основной теореме алгебры любой многочлен над $\C$ имеет корень:
	\[
	f(x)=(x-\lambda)g(x).
	\]
	Следовательно, неприводимы только линейные многочлены.
\end{proof}

\begin{theorem}
	Неприводимые многочлены над $\R$ имеют степень $1$ или $2$.
\end{theorem}

\begin{proof}
	Комплексные корни вещественного многочлена возникают попарно:
	\[
	a\pm bi.
	\]
	Поэтому неприводимые квадратичные множители имеют вид
	\[
	x^2+px+q,
	\]
	где дискриминант отрицателен.
\end{proof}

\begin{theorem}
	Над любым полем существует бесконечно много унитарных неприводимых
	многочленов.
\end{theorem}

\begin{proof}
	Предположим противное:
	\[
	f_1,\dots,f_n
	\]
	--- все неприводимые унитарные многочлены.
	
	Рассмотрим
	\[
	F=f_1f_2\cdots f_n+1.
	\]
	Многочлен $F$ не делится ни на один из $f_i$, противоречие.
\end{proof}

\begin{theorem}
	Над конечным полем существуют неприводимые многочлены сколь угодно большой степени.
\end{theorem}

\begin{proof}
	Если бы степени были ограничены числом $N$, то существовало бы лишь конечное
	число неприводимых многочленов. Но тогда и всех многочленов было бы конечное
	число, противоречие.
\end{proof}

%──────────────────────────────────────────────────────────────────────────────
\section{Примитивные многочлены и лемма Гаусса}

\begin{definition}
	Многочлен над факториальным кольцом называется примитивным, если НОД его
	коэффициентов равен $1$.
\end{definition}

\begin{theorem}
	Любой многочлен
	\[
	f(x)\in K[x]
	\]
	можно представить в виде
	\[
	f(x)=c\cdot g(x),
	\]
	где
	\[
	c\in K,
	\]
	а $g(x)$ --- примитивный многочлен.
\end{theorem}

\begin{theorem}[Лемма Гаусса]
	Произведение примитивных многочленов является примитивным.
\end{theorem}

\begin{proof}
	Пусть
	\[
	f,g
	\]
	--- примитивные многочлены. Предположим, что
	\[
	fg
	\]
	непримитивен. Тогда существует простой элемент $p$, делящий все коэффициенты
	произведения.
	
	Возьмём первые коэффициенты в $f$ и $g$, не делящиеся на $p$. Их произведение
	даёт коэффициент в $fg$, не делящийся на $p$, противоречие.
\end{proof}

%──────────────────────────────────────────────────────────────────────────────
\section{Неприводимость над $\Z$ и $\Q$}

\begin{theorem}
	Многочлен из $\Z[x]$ примитивен и неприводим над $\Z$ тогда и только тогда,
	когда он неприводим над $\Q$.
\end{theorem}

\begin{proof}
	Если многочлен раскладывается над $\Q$, после приведения к общему знаменателю
	и применения леммы Гаусса он раскладывается и над $\Z$.
\end{proof}

\begin{theorem}[Критерий Эйзенштейна]
	Пусть существует простое число $p$ такое, что:
	\[
	p\mid a_0,\dots,a_{n-1},
	\]
	\[
	p\nmid a_n,
	\]
	\[
	p^2\nmid a_0.
	\]
	Тогда многочлен
	\[
	a_nx^n+\dots+a_0
	\]
	неприводим над $\Q$.
\end{theorem}

\begin{theorem}[Редукционный критерий]
	Если многочлен из $\Z[x]$ после редукции по модулю простого числа $p$
	остаётся неприводимым, то он неприводим над $\Q$.
\end{theorem}

%──────────────────────────────────────────────────────────────────────────────
\section{Факториальность кольца многочленов}

\begin{theorem}
	Если кольцо $R$ факториально, то кольцо
	\[
	R[x]
	\]
	тоже факториально.
\end{theorem}

\begin{proof}
	Каждый многочлен раскладывается на примитивную часть и содержимое.
	Неприводимость изучается в поле частных кольца $R$.
	
	С помощью леммы Гаусса разложение в поле переносится в $R[x]$, а
	единственность сохраняется.
\end{proof}

%──────────────────────────────────────────────────────────────────────────────
\section{Многочлены от нескольких переменных. Степень, однородные многочлены. Область целостности и факториальность кольца многочленов. Существование неприводимых многочленов любой степени над любым полем}
\addcontentsline{toc}{section}{Вопрос 60. Многочлены от нескольких переменных. Степень, однородные многочлены. Область целостности и факториальность кольца многочленов. Существование неприводимых многочленов любой степени над любым полем}

Кольцо многочленов от $n$ переменных определяется индуктивно: $R[x_1, \dots, x_n] := (R[x_1, \dots, x_{n-1}])[x_n]$. Элементами этого кольца являются конечные суммы формальных выражений.

\begin{definition}
	\textbf{Многочленом} от $n$ переменных $x_1, \dots, x_n$ над коммутативным кольцом $R$ называется формальная сумма вида:
	\[
	f(x_1, \dots, x_n) = \sum_{(i_1, \dots, i_n)} a_{i_1 \dots i_n} x_1^{i_1} x_2^{i_2} \cdots x_n^{i_n},
	\]
	где $a_{i_1 \dots i_n} \in R$, а суммирование ведется по конечному множеству наборов неотрицательных целых чисел $(i_1, \dots, i_n) \in \mathbb{Z}_{\ge 0}^n$. Каждое отдельное слагаемое $a_{\mathbf{i}} x_1^{i_1} \cdots x_n^{i_n}$ называется \textbf{одночленом} (мономом).
\end{definition}

\begin{definition}
	\textbf{Степенью одночлена} (тотальной степенью) $x_1^{i_1} \cdots x_n^{i_n}$ называется сумма его показателей: $\deg(x_1^{i_1} \cdots x_n^{i_n}) = i_1 + i_2 + \dots + i_n$. Степенью всего многочлена $f \neq 0$ называется максимум степеней входящих в него одночленов с ненулевыми коэффициентами.
\end{definition}

\begin{definition}
	Многочлен $f(x_1, \dots, x_n)$ называется \textbf{однородным} (или формой степени $d$), если абсолютно все входящие в него одночлены с ненулевыми коэффициентами имеют одну и ту же степень $d$.
\end{definition}

\begin{remark}
	Для корректного определения старшего члена вводится \textbf{лексикографический порядок} на наборах степеней. Мы говорим, что $(i_1, \dots, i_n) > (j_1, \dots, j_n)$, если для первого индекса $k$, где $i_k \neq j_k$, выполнено $i_k > j_k$. Одночлен с ненулевым коэффициентом, имеющий максимальный в смысле этого порядка набор экспонент, называется \textbf{старшим одночленом} многочлена.
\end{remark}

\begin{theorem}
	Если $R$ — область целостности, то кольцо многочленов $R[x_1, \dots, x_n]$ также является областью целостности.
\end{theorem}
\begin{proof}
	Применим математическую индукцию по числу переменных $n$.
	\begin{itemize}
		\item \textbf{База индукции ($n=1$):} Пусть $f(x), g(x) \in R[x]$ — два ненулевых многочлена. Пусть их старшие члены равны $a_k x^k$ ($a_k \neq 0$) и $b_m x^m$ ($b_m \neq 0$) соответственно. Тогда старший член их произведения $f(x)g(x)$ равен $(a_k b_m) x^{k+m}$. Поскольку $R$ — область целостности, произведение ненулевых элементов $a_k b_m \neq 0$. Следовательно, $f(x)g(x) \neq 0$, и делителей нуля в $R[x]$ нет.
		\item \textbf{Индукционный переход:} Предположим, что утверждение верно для $n-1$ переменных, то есть $S = R[x_1, \dots, x_{n-1}]$ — область целостности. Но по определению $R[x_1, \dots, x_n] \cong S[x_n]$. По базе индукции, кольцо многочленов от одной переменной $x_n$ над областью целостности $S$ само является областью целостности. Переход доказан.
	\end{itemize}
\end{proof}

\begin{theorem}
	Если кольцо $R$ факториально (является ООР), то кольцо многочленов $R[x_1, \dots, x_n]$ также факториально.
\end{theorem}
\begin{proof}
	Доказательство опирается на \textbf{теорему Гаусса}, которая гласит: \textit{«Если кольцо $A$ факториально, то и кольцо $A[x]$ факториально»}.
	Снова применим индукцию по числу переменных $n$:
	\begin{itemize}
		\item \textbf{База ($n=1$):} Прямо следует из теоремы Гаусса для кольца $R[x_1]$.
		\item \textbf{Переход от $n-1$ к $n$:} Пусть кольцо $S = R[x_1, \dots, x_{n-1}]$ факториально по предположению индукции. Рассматривая $R[x_1, \dots, x_n]$ как $S[x_n]$ и применяя теорему Гаусса к факториальному кольцу коэффициентов $S$, получаем, что $S[x_n]$ факториально.
	\end{itemize}
\end{proof}

\begin{theorem}
	Над любым полем $K$ существуют неприводимые многочлены любой заданной степени $d \ge 1$.
\end{theorem}
\begin{proof}
	Разберем два случая в зависимости от природы поля $K$:
	\begin{enumerate}
		\item \textbf{Случай 1. Поле $K$ — бесконечное.} 
		Для степени $d=1$ любой многочлен вида $x - a$ неприводим. Для степеней $d > 1$ существование гарантируется, например, теоремой Кронеккера и теоремой о примитивном элементе. В частности, если взять поле рациональных чисел $\mathbb{Q}$, то по критерию Эйзенштейна многочлен $x^d - 2$ всегда неприводим над $\mathbb{Q}$ для любой степени $d$.
		
		\item \textbf{Случай 2. Поле $K = \mathbb{F}_q$ — конечное ($q = p^m$).}
		Из теории конечных полей известно, что существует единственное расширение поля $\mathbb{F}_q$ степени $d$, а именно поле $\mathbb{F}_{q^d}$. Мультипликативная группа этого поля циклическая, и поле может быть порождено одним примитивным элементом $\alpha$, то есть $\mathbb{F}_{q^d} = \mathbb{F}_q(\alpha)$. 
		
		Минимальный многочлен $m_\alpha(x) \in \mathbb{F}_q[x]$ этого элемента $\alpha$ обязан быть:
		а) неприводимым над $\mathbb{F}_q$ по определению минимального многочлена;
		б) его степень $\deg(m_\alpha) = [\mathbb{F}_q(\alpha) : \mathbb{F}_q] = d$.
		Таким образом, минимальный многочлен элемента, порождающего расширение степени $d$, и является искомым неприводимым многочленом степени $d$.
	\end{enumerate}
\end{proof}

%──────────────────────────────────────────────────────────────────────────────
\section{Доказать, что $R[x,y,\dots,z]$, где $R$ — поле, не является кольцом главных идеалов}
\addcontentsline{toc}{section}{Вопрос 61. Доказать, что $R[x,y,\dots,z]$, где $R$ — поле, не является кольцом главных идеалов}

\begin{theorem}
	Кольцо многочленов от двух и более переменных $K[x, y]$ над полем $K$ не является кольцом главных идеалов.
\end{theorem}
\begin{proof}
	Рассмотрим идеал $I = \langle x, y \rangle$, порожденный переменными $x$ и $y$. Этот идеал состоит из всех многочленов без свободного члена. 
	Предположим, что $I$ — главный, то есть существует такой многочлен $d(x,y)$, что $I = \langle d(x,y) \rangle$.
	Тогда $x \in \langle d \rangle \Rightarrow d \mid x$. Отсюда следует, что степень $d$ по переменной $y$ равна 0, а по $x$ равна либо 0, либо 1.
	Аналогично, $y \in \langle d \rangle \Rightarrow d \mid y$, откуда степень $d$ по $x$ равна 0.
	Таким образом, $d(x,y) = c \in K^{\times}$ — константа. Но если $d$ константа, то $\langle d \rangle = K[x,y]$. Однако наш идеал $I \neq K[x,y]$, так как $1 \notin I$. Получили противоречие. Идеал не является главным.
\end{proof}

%──────────────────────────────────────────────────────────────────────────────
\section{Симметрические многочлены. Алгебраически зависимые элементы}
\addcontentsline{toc}{section}{Вопрос 62. Симметрические многочлены. Алгебраически зависимые элементы}

\begin{definition}
	Многочлен $f(x_1, \dots, x_n)$ называется \textbf{симметрическим}, если он не меняется при любой перестановке своих переменных.
\end{definition}

\begin{definition}
	Элементарные симметрические многочлены задаются как коэффициенты многочлена по переменной $t$ с корнями $x_1, \dots, x_n$:
	\begin{align*}
		\sigma_1 &= x_1 + x_2 + \dots + x_n, \\
		\sigma_2 &= \sum_{i < j} x_i x_j, \\
		&\dots \\
		\sigma_n &= x_1 x_2 \dots x_n.
	\end{align*}
\end{definition}

\begin{theorem}[Основная теорема о симметрических многочленах]
	Любой симметрический многочлен единственным образом выражается в виде многочлена от элементарных симметрических многочленов.
\end{theorem}
\begin{proof}
	Доказательство проводится по индукции с использованием лексикографического порядка старших членов. Из исходного симметрического многочлена последовательно вычитаются подходящие комбинации $\sigma_1^{k_1}\dots\sigma_n^{k_n}$, уничтожающие текущий старший член. Процесс завершается за конечное число шагов.
\end{proof}

%──────────────────────────────────────────────────────────────────────────────
\section{Теорема Кронеккера}

\begin{theorem}[Теорема Кронеккера]
	Для любого многочлена
	\[
	f(x)\in K[x]
	\]
	существует расширение поля, в котором этот многочлен раскладывается на
	линейные множители.
\end{theorem}

\begin{proof}
	Если многочлен неприводим, присоединим к полю его корень:
	\[
	K[x]/(f).
	\]
	Тогда
	\[
	f(x)=(x-\alpha)g(x).
	\]
	Продолжая процесс для многочлена $g(x)$, после конечного числа шагов
	получим поле разложения.
\end{proof}

%──────────────────────────────────────────────────────────────────────────────
\section{Общие корни и результант}

\begin{definition}
	Результантом многочленов
	\[
	f,g
	\]
	называется определитель матрицы Сильвестра:
	\[
	\operatorname{Res}(f,g).
	\]
\end{definition}

\begin{theorem}
	Многочлены имеют общий корень тогда и только тогда, когда
	\[
	\operatorname{Res}(f,g)=0.
	\]
\end{theorem}

\begin{proof}
	Если многочлены имеют общий корень $\alpha$, то
	\[
	f(\alpha)=g(\alpha)=0,
	\]
	поэтому строки матрицы Сильвестра линейно зависимы.
	
	Обратно, если результант равен нулю, существуют многочлены
	\[
	u,v
	\]
	такие, что
	\[
	uf+vg=0.
	\]
	Следовательно, $f$ и $g$ имеют общий неприводимый делитель, а значит,
	общий корень в некотором расширении поля.
\end{proof}

\end{document}
